% !TEX root = ../main.tex
\chapter{谱理论与迹公式}\label{ch:08}

\begin{center}
Daniel Bump\footnote{作者感谢 Yonatan Gutman 与审稿人对初稿的详尽评论与诸多帮助。本章的写作部分得到美国国家科学基金 DMS-9970841 的资助。}
\end{center}

我们叙述 $\Gamma\backslash\mathrm{SL}(2,\mathbb{R})$ 的谱理论的一部分，特别是 Selberg 迹公式，
并侧重表示论中的想法。为简单起见，迹公式只在紧商的情形下处理。
最后一节性质不同，意在展示迹公式在提升问题上的一项简单应用。

\section{谱问题}

群 $G=\mathrm{SL}(2,\mathbb{R})$ 作用在 $\mathfrak{H}=\{x+iy\mid y>0\}$ 上：
\[
\begin{pmatrix}a&b\\c&d\end{pmatrix}:z\longmapsto\frac{az+b}{cz+d}.
\]
$i$ 的稳定化子是 $K=\mathrm{SO}(2)$，故 $\mathfrak{H}\cong G/K$。\emph{非欧 Laplacian}
\begin{equation}\tag{1}
\Delta=-y^{2}\left(\frac{\partial^{2}}{\partial x^{2}}+\frac{\partial^{2}}{\partial y^{2}}\right)
\end{equation}
是 $G$-不变微分算子。设 $\Gamma$ 是 $G$ 的离散余紧子群，则 $X=\Gamma\backslash\mathfrak{H}$
是紧 Riemann 曲面。当前的问题就是描述 $\Delta$ 在 $X$ 上的谱。

\par\medskip\noindent\textbf{命题 1.}\quad\emph{群 $\Gamma$ 在 $\mathfrak{H}$ 上有基本区域
$\mathcal{F}$，其边界由成对的测地弧 $\alpha_{i}$ 与 $\gamma_{i}(\alpha_{i})$ 组成，其中
$\gamma_{i}\in\Gamma$。当边界按逆时针方向走时，全等的弧 $\alpha_{i}$ 与 $\gamma_{i}(\alpha_{i})$
的走向相反。}

\smallskip
\noindent\emph{证明梗概.}\ 取点 $P\in\mathfrak{H}$，它除 $\pm I$ 外不被 $\Gamma$ 的任何元素固定。
令 $\mathcal{F}$ 为距 $P$ 比距 $\gamma(P)$ 更近的点集，距离按非欧度量计算，并对所有
$\gamma\in\Gamma$ 取交。令 $N=\{\gamma_{1},\dots,\gamma_{n}\}$ 为使 $\gamma_{i}(\mathcal{F})$
与 $\mathcal{F}$ 相邻的 $\Gamma$ 的元素之集。

先设没有 $\gamma_{i}^{2}=1$。显然每个 $\gamma_{i}^{-1}\in N$，故 $N$ 有偶数个元素 $2h$。
重排下标使 $\gamma_{i+h}=\gamma_{i}^{-1}$（$1\le i\le h$）。令 $\alpha_{i}$ 为下述点集的测地线
与 $\mathcal{F}$ 的闭包之交：该点集由到 $P$ 与到 $\gamma_{i}^{-1}(P)$ 等距的点组成。
于是 $\gamma_{i}(\alpha_{i})=\alpha_{i+h}$，而 $\alpha_{1},\dots,\alpha_{h}$ 满足要求。

若某个 $\gamma_{i}^{2}=1$，则弧 $\alpha_{i}$ 在 $\gamma_{i}$ 下自全等，其中点被 $\gamma_{i}$ 固定。
把这种弧在中点一分为二，结论同样成立。\hfill$\square$

\medskip
Laplacian $\Delta$ 作用在 $C^{\infty}(\Gamma\backslash\mathfrak{H})$ 上。

\par\medskip\noindent\textbf{命题 2.}\quad\emph{$\Delta$ 关于不变度量 $y^{-2}dx\wedge dy$
是对称的、正的。}

\smallskip
\noindent\emph{证明.}\ 对称性指
\[
\langle\Delta f,g\rangle-\langle f,\Delta g\rangle=0,\qquad f,g\in C^{\infty}(\Gamma\backslash\mathfrak{H}).
\]
取命题 1 中的基本区域 $\mathcal{F}$，上式等于
\[
\int_{\mathcal{F}}\left(f\left(\frac{\partial^{2}\bar g}{\partial x^{2}}+\frac{\partial^{2}\bar g}{\partial y^{2}}\right)
-\bar g\left(\frac{\partial^{2}f}{\partial x^{2}}+\frac{\partial^{2}f}{\partial y^{2}}\right)\right)dx\wedge dy
=\int_{\mathcal{F}}d\omega,
\]
其中
\[
\omega=f\,\frac{\partial\bar g}{\partial x}\,dy-f\,\frac{\partial\bar g}{\partial y}\,dx
-\bar g\,\frac{\partial f}{\partial x}\,dy+\bar g\,\frac{\partial f}{\partial y}\,dx.
\]
由 Stokes 定理，此积分等于 $\omega$ 沿 $\mathcal{F}$ 的边界的积分。由命题 1，相邻边界弧的贡献
两两相消。

正性指 $\langle\Delta f,f\rangle\ge0$，且等号仅在 $f$ 为常数时成立。我们计算
\[
\left[\left|\frac{\partial f}{\partial x}\right|^{2}+\left|\frac{\partial f}{\partial y}\right|^{2}
+\left(\frac{\partial^{2}f}{\partial x^{2}}+\frac{\partial^{2}f}{\partial y^{2}}\right)\bar f\right]dx\wedge dy
=d\left(\bar f\left(\frac{\partial f}{\partial x}\,dy-\frac{\partial f}{\partial y}\,dx\right)\right).
\]
再用 Stokes 定理，即得
\begin{equation}\tag{2}
\langle\Delta f,f\rangle\ge-\int_{\partial\mathcal{F}}\bar f\left(\frac{\partial f}{\partial x}\,dy-\frac{\partial f}{\partial y}\,dx\right),
\end{equation}
等号仅在 $\partial f/\partial x=\partial f/\partial y=0$ 恒成立、即 $f$ 为常数时取到。
由 Cauchy--Riemann 方程与链式法则可以验证：若 $x+iy$ 与 $u+iv$ 由全纯映射相联，
例如由线性分式变换相联，则
\[
\frac{\partial f}{\partial x}\,dy-\frac{\partial f}{\partial y}\,dx
=\frac{\partial f}{\partial u}\,dv-\frac{\partial f}{\partial v}\,du.
\]
由此可见，全等边界弧的贡献相消，(2) 的右端为零。\hfill$\square$

后面我们会看到，$\Delta$ 延拓为 $\mathcal{L}=L^{2}(\Gamma\backslash\mathfrak{H},y^{-2}dx\wedge dy)$
上的自伴无界算子。这指的是：存在 $\mathcal{L}$ 的稠密子空间 $D$，它包含
$C^{\infty}(\Gamma\backslash\mathfrak{H})$，使 $\Delta$ 延拓到 $D$；并且若 $v\in\mathcal{L}$ 使
$u\mapsto\langle v,\Delta u\rangle$ 在 $D$ 上连续，则 $v\in D$，且
$\langle\Delta v,u\rangle=\langle v,\Delta u\rangle$。

$\Delta$ 自伴，故有一套良好的谱理论，这正是我们要发展的。
具体做法是引入与 $\Delta$ 交换的积分算子，证明它们迹类，并对它们证明 Selberg 迹公式。

最后我们考虑 $\Gamma$ 带一个尖点的较难情形。此时离散谱与连续谱并存，
而 Eisenstein 级数的理论是这一情形的本质特征。

\begin{annote}
原书取 $\Delta=-y^{2}(\partial_{x}^{2}+\partial_{y}^{2})$，这是上半平面双曲度量
$y^{-2}(dx^{2}+dy^{2})$ 下 Laplace--Beltrami 算子的正值形式，其谱含于 $[0,\infty)$。
特征值通常写成 $\lambda=s(1-s)$，即 $s=\tfrac12\pm it$：当 $\lambda\ge\tfrac14$ 时 $t$ 为实数，
对应主系列；当 $0<\lambda<\tfrac14$ 时 $t$ 为纯虚数，对应补系列；$\lambda=0$ 对应 $s=0,1$。
但第 3 节为方便起见令 $\lambda=-s$ 且 $s>0$，那里的 $s$ 与本节之后 $s=\tfrac12+it$ 的 $s$
并非同一参数，读 (7) 与定理 18 时务必分开。
\par\vspace{0.15em}{\raggedleft 【来源：Hejhal, \emph{The Selberg Trace Formula for PSL(2,R)}, Vol.~2, Springer LNM 1001, 1983, Ch.~1；Iwaniec, \emph{Introduction to the Spectral Theory of Automorphic Forms}, Madrid, 1995, Ch.~1（A）】\par}
\end{annote}

\section{积分算子环}

群 $G$ 通过\emph{右正则作用} $(\rho(g)f)(x)=f(xg)$ 作用在函数 $f:G\to\mathbb{C}$ 上。
既然马上要讨论微分算子，我们暂且把自己限制在 $f\in C^{\infty}(G)$。

由于 $\mathfrak{H}$ 可与齐性空间 $G/K$ 典范地等同，$\mathfrak{H}$ 上的函数 $f$ 可以看作 $G$
上右 $K$-不变的函数。$G$ 的正则表示与 $\mathfrak{H}$ 上函数的关系并不显然：
右平移并不保持右 $K$-不变性。但两者确有联系，下面说明。

$G$ 的李代数是 $\mathfrak{g}$，即迹为零的 $2\times2$ 实矩阵构成的空间。
它通过
\[
Xf(g)=\frac{d}{dt}f(ge^{tX})\Big|_{t=0},\qquad X\in\mathfrak{g},
\]
作用在 $G$ 的光滑函数上。这是李代数表示：$[X,Y]=XY-YX$ 的作用等于 $X\circ Y-Y\circ X$，
其中 $X\circ Y$ 是算子的复合。这个 $\mathfrak{g}$ 的表示即正则表示的微分形式。

泛包络代数 $U(\mathfrak{g})$ 是由 $\mathfrak{g}$ 生成、并满足关系
\[
[X,Y]-(X\cdot Y-Y\cdot X)=0
\]
的结合 $\mathbb{R}$-代数，其中 $X\cdot Y$ 表示 $U(\mathfrak{g})$ 中的乘法。
任何李代数表示都唯一地延拓为 $U(\mathfrak{g})$ 的表示，特别地，$\mathfrak{g}$ 的正则表示延拓为
$U(\mathfrak{g})$ 的表示。于是 $U(\mathfrak{g})$ 实现为 $C^{\infty}(G)$ 上的一个微分算子环。

若 $D$ 是 $U(\mathfrak{g})$ 的\emph{中心}的元素，则它与正则表示交换。这一点符合直觉，
Bump [6, Proposition 2.2.4] 中有证明。特别地，若 $f$ 在 $\rho(K)$ 下不变，则 $Df$ 也不变，
因此 $D$ 作用在 $C^{\infty}(\mathfrak{H})$ 上。

$U(\mathfrak{g})$ 的中心的一个特定元素是 \emph{Laplace--Beltrami} 元或 Casimir 元
\begin{equation}\tag{3}
-4\Delta=H\cdot H+2R\cdot L+2L\cdot R,
\end{equation}
其中
\[
R=\begin{pmatrix}0&1\\0&0\end{pmatrix},\qquad
L=\begin{pmatrix}0&0\\1&0\end{pmatrix},\qquad
H=\begin{pmatrix}1&0\\0&-1\end{pmatrix}.
\]
事实上 $\mathbb{C}[\Delta]$ 就是 $U(\mathfrak{g})$ 的中心。我们在此沿用了先前表示非欧 Laplacian
的字母 $\Delta$，因为当 $U(\mathfrak{g})$ 的这个元素解释为 $\mathfrak{H}$ 上的微分算子时，
两者相同。见 Bump [6, Proposition 2.2.5]。

令 $\mathcal{H}$ 为 $G$ 上光滑、紧支函数关于卷积构成的环，$\mathcal{H}^{\circ}$ 为其中
$K$-双不变函数构成的子环。两者都没有单位元。我们把 $\mathcal{H}^{\circ}$ 称为
\emph{球 Hecke 代数}，但要提醒读者：另有几个自然且密切相关的环也曾用这个名称。
环 $\mathcal{H}$ 非交换，但我们有下述结论。

\par\medskip\noindent\textbf{定理 3.}\quad\emph{环 $\mathcal{H}^{\circ}$ 交换。}

\smallskip
\noindent\emph{证明.}\ 矩阵转置保持 $K$，故在 $\mathcal{H}^{\circ}$ 上诱导对合 $\iota$，满足
$\iota(\phi*\psi)=\iota(\psi)*\iota(\phi)$，其中 $(\iota f)(g)=f(g^{t})$。
每个双陪集 $K\backslash G/K$ 都有对角代表元，故 $\iota$ 是恒等映射。\hfill$\square$

Gelfand 的这个定理有表示论的含义。设 $(\pi,V)$ 是 $G$ 在 Banach 空间上的表示，
以 $V^{K}$ 记 $K$-不变向量构成的子空间。代数 $\mathcal{H}$ 通过
\[
\pi(\phi)v=\int_{G}\phi(g)\pi(g)v\,dg,\qquad \phi\in\mathcal{H},\ v\in V
\]
作用在 $V$ 上。这里 $\int_{G}dg$ 是 Haar 积分。若 $\phi\in\mathcal{H}^{\circ}$，则
$\pi(\phi)v\in V^{K}$。特别地，$V^{K}$ 是 $\mathcal{H}^{\circ}$ 的模。

$K$ 紧，故 $K$ 的任何不可约表示 $\rho$ 都有限维。$V$ 的 $\rho$-等型部分 $V(\rho)$
定义为所有与 $\rho$ 作为 $K$-模同构的 $K$-不变子空间的直和。$G$ 的表示 $(\pi,V)$
称为\emph{容许的}，若每个 $V(\rho)$ 都有限维。特别地，取 $\rho$ 为平凡表示，
此时 $V(\rho)=V^{K}$，故 $V^{K}$ 必须有限维。

\par\medskip\noindent\textbf{定理 4.}\quad\emph{若 $(\pi,V)$ 是 $G$ 的不可约容许表示，
则 $\dim V^{K}\le1$。}

\smallskip
\noindent\emph{证明.}\ $\mathcal{H}^{\circ}$ 交换，故其有限维不可约模都是一维的。
于是只需证明 $V^{K}$（若非零）是 $\mathcal{H}^{\circ}$ 的不可约模。
设 $L\subset V^{K}$ 是非零闭的 $\mathcal{H}^{\circ}$-不变子空间。若 $v\in V^{K}$，
我们要证 $v\in L$。

任给 $\epsilon>0$。$V$ 不可约而 $L$ 非零，故 $\pi(\mathcal{H})L$ 的闭包是 $V$，
于是存在 $\phi\in\mathcal{H}$、$w\in L$ 使 $\pi(\phi)w=v_{1}$，其中 $|v_{1}-v|<\epsilon$。
令 $v_{2}=\int_{K}\pi(k)v_{1}\,dk$。则 $v_{2}$ 是 $K$-不变的，而 $v$ 也 $K$-不变，故
\[
|v_{2}-v|=\left|\int_{K}\pi(k)(v_{1}-v)\,dk\right|
\le\int_{K}|\pi(k)(v_{1}-v)|\,dk=\int_{K}|v_{1}-v|\,dk<\epsilon .
\]
这里把 Haar 测度归一化使 $K$ 的体积为 $1$。由于 $w$ 是 $K$-不变的，
\[
\pi(k)\pi(\phi)\pi(k')w=\pi(k)v_{1}
\]
对一切 $k,k'\in K$ 成立。对 $k$ 与 $k'$ 积分，即得 $\pi(\phi_{0})w=v_{2}$，其中
\[
\phi_{0}(g)=\int_{K\times K}\phi(kgk')\,dk\,dk'.
\]
现在 $\phi_{0}\in\mathcal{H}^{\circ}$，而 $L$ 是 $\mathcal{H}^{\circ}$-不变的，故 $v_{2}\in L$。
于是 $v$ 可被 $L$ 的元素任意逼近；又 $L\subset V^{K}$ 有限维、因而闭，所以 $v\in L$。
\hfill$\square$

更一般地，若 $(\pi,V)$ 是 $G$ 的表示、$k\in\mathbb{Z}$，我们用 $V(k)$ 记 $V$ 中满足
\[
\pi(\kappa_{\theta})v=e^{ik\theta}v,\qquad
\kappa_{\theta}=\begin{pmatrix}\cos\theta&\sin\theta\\-\sin\theta&\cos\theta\end{pmatrix}\in K
\]
的 $v$ 构成的子空间。这就是 $\rho$-等型子空间，其中 $\rho$ 是 $K$ 的特征 $e^{ik\theta}$。
若 $(\pi,V)$ 不可约，则因 $-I\in K$ 是中心元，它作用为纯量 $(-1)^{\epsilon}$，
$\epsilon=0$ 或 $1$。显然，要使 $V(k)$ 非零，$k$ 的奇偶性必须与 $\epsilon$ 一致。
$\epsilon=0$ 时称 $\pi$ 为\emph{偶}的，$\epsilon=1$ 时称为\emph{奇}的。

\par\medskip\noindent\textbf{命题 5.}\quad\emph{若 $(\pi,V)$ 是 $G$ 的不可约容许表示，
则 $\dim V(k)\le1$。}

\smallskip
\noindent\emph{证明梗概.}\ 可沿定理 4 的同一思路证明。由于 $K=\mathrm{SO}(2)$ 交换，
可以验证：满足 $f(\kappa_{\theta}g\kappa_{\theta'})=e^{ik(\theta+\theta')}f(g)$ 的光滑紧支函数
所成的卷积环交换。见 Bump [6, Proposition 2.2.8]。\hfill$\square$

我们指出，定理 4 可直接推广到任意约化李群，命题 5 则不能。于是对约化李群的不可约容许表示，
其极大紧子群平凡表示的重数至多为 $1$；极大紧子群的其他不可约表示各以有限重数出现，
这就是容许性，但重数不必为 $1$。事实上不可约表示自动容许，不过自守形式理论并不需要这一事实，
那里自守表示的容许性可直接证明。

若 $(\pi,V)$ 是 $G$ 的不可约表示，以 $V_{\mathrm{fin}}$ 记诸 $V(k)$ 的直和，
这就是 $K$-有限向量构成的空间。它不在 $G$ 的作用下不变，但在李代数 $\mathfrak{g}$ 与 $K$
的作用下不变，因此是一个 $(\mathfrak{g},K)$-模。关于这个概念，见 Bump [6, p.\ 200]。

若 $\pi=\rho$ 是 $G$ 的正则表示，我们不用 $\rho(\phi)$ 而用 $T_{\phi}$ 记相应的算子，于是
\[
T_{\phi}f(x)=\int_{G}\phi(g)f(xg)\,dg .
\]
这是与 $g\mapsto\phi(g^{-1})$ 的卷积。记号意在提示 Hecke 算子。
这些积分算子对我们很重要，因为它们与 $\Delta$ 交换，却比 $\Delta$ 容易研究。
注意 $L^{2}(\Gamma\backslash G)$ 在 $T_{\phi}$ 下不变；若 $\phi\in\mathcal{H}^{\circ}$，
则 $T_{\phi}$ 保持右 $K$-不变性，因而可以看作 $L^{2}(\Gamma\backslash\mathfrak{H})$ 上的积分算子。

设 $\phi\in\mathcal{H}$，置
\begin{equation}\tag{4}
K_{\phi}(x,y)=\sum_{\gamma\in\Gamma}\phi(x^{-1}\gamma y).
\end{equation}
起初我们把 $(x,y)$ 看作 $G\times G$ 的元素。若 $x$ 与 $y$ 限制在紧集 $C$ 上，
则 $\phi(x^{-1}\gamma y)$ 当 $\gamma$ 落在某个紧集 $C'$ 之外时为零，故只有有限多个 $\gamma$
有贡献。于是 (4) 收敛，并定义 $x$ 与 $y$ 的光滑函数。

变量替换表明，$K_{\phi}(x,y)$ 在 $x$ 或 $y$ 左乘 $\gamma$ 时不变，因此这个核既可看作
$G\times G$ 上的函数，也可看作 $\Gamma\backslash G\times\Gamma\backslash G$ 上的函数，
且连续。若 $\phi\in\mathcal{H}^{\circ}$，则 $K_{\phi}$ 甚至可以看作 $X\times X$ 上的函数，
其中 $X=\Gamma\backslash\mathfrak{H}$。

\par\medskip\noindent\textbf{命题 6.}\quad\emph{有}
\begin{equation}\tag{5}
(T_{\phi}f)(x)=\int_{\Gamma\backslash G}K_{\phi}(x,y)\,f(y)\,dy .
\end{equation}

\smallskip
\noindent\emph{证明.}\ 左端等于
\[
\int_{G}\phi(y)f(xy)\,dy=\int_{G}\phi(x^{-1}y)f(y)\,dy
=\sum_{\gamma\in\Gamma}\int_{\Gamma\backslash G}\phi(x^{-1}\gamma y)f(y)\,dy,
\]
其中用到 $f(\gamma y)=f(y)$。交换求和与积分即得 (5)。\hfill$\square$

设 $H$ 是 Hilbert 空间。算子 $T:H\to H$ 称为\emph{紧的}，若 $T$ 把有界集映为相对紧集。

\par\medskip\noindent\textbf{定理 7.}\quad\emph{$T_{\phi}$ 是 $L^{2}(\Gamma\backslash G)$
上的紧算子。}

\smallskip
见 Bump [6, Section II.3]，特别是 Theorem 2.3.2 与 Proposition 2.3.1，那里有更完整的细节。

\smallskip
\noindent\emph{证明.}\ 核 $K_{\phi}$ 在紧空间 $(\Gamma\backslash G)\times(\Gamma\backslash G)$
上连续，故必属于 $L^{2}((\Gamma\backslash G)\times(\Gamma\backslash G))$。
Hilbert 与 Schmidt 的著名定理断言：若 $Z$ 是局部紧 Borel 测度空间、$L^{2}(Z)$ 可分，
则核为 $K\in L^{2}(Z\times Z)$ 的积分算子
\[
(Tf)(x)=\int_{Z}K(x,y)f(y)\,dy
\]
是紧的。\hfill$\square$

若
\begin{equation}\tag{6}
\phi(g^{-1})=\overline{\phi(g)},
\end{equation}
则 $K_{\phi}(x,y)=\overline{K_{\phi}(y,x)}$，故 $T_{\phi}$ 自伴。

\par\medskip\noindent\textbf{定理 8.}\quad\emph{设 $T$ 是可分 Hilbert 空间 $H$ 上的紧自伴算子。
则 $H$ 有一组由 $T$ 的特征向量构成的标准正交基 $\phi_{i}$（$i=1,2,3,\dots$），
使 $T\phi_{i}=\mu_{i}\phi_{i}$；特征值满足 $\mu_{i}\to0$（$i\to\infty$）。}

\smallskip
这就是紧算子的谱定理。见 Bump [6, Theorem 2.3.1]。\hfill$\square$

于是若 (6) 成立，$T_{\phi}$ 是自伴紧算子，其非零特征值 $\mu_{i}\to0$。
Hilbert--Schmidt 性质给出更多信息：$\sum|\mu_{i}|^{2}<\infty$。
后面我们会看到更多：$\sum|\mu_{i}|<\infty$。这意味着 $T_{\phi}$ 是迹类算子。
这一事实很重要，因为 Selberg 迹公式正是关于它的。

\par\medskip\noindent\textbf{定理 9.}\quad\emph{$L^{2}(X)$ 有一组由 $\Delta$ 的特征向量
构成的基。}

\smallskip
\noindent\emph{证明.}\ 满足 (6) 的诸 $T_{\phi}$ 是一族彼此交换的自伴紧算子，
故可同时对角化。由谱定理，非零特征子空间有限维；又不存在使所有 $T_{\phi}$ 都为零的非零向量，
因为可以取 $\phi$ 为正、总量为 $1$ 且集中在单位元附近，此时 $T_{\phi}f$ 逼近 $f$。
因此 $\mathcal{H}^{\circ}$ 的各同时特征子空间都有限维。

设 $V$ 是这样一个特征子空间。$\Delta$ 与诸 $T_{\phi}$ 交换，故保持 $V$；
又 $\Delta$ 对称，故在 $V$ 上诱导自伴变换。在每个这样的 $V$ 中取一组由 $\Delta$
的特征向量构成的标准正交基，再把所有 $V$ 的基并起来即可。\hfill$\square$

与定理 9 密切相关的是关于 $L^{2}(\Gamma\backslash G)$ 的一个表示论命题。
正则表示 $\rho$ 是这个空间上的酉表示。

\par\medskip\noindent\textbf{引理 10.}\quad\emph{设 $H$ 是 $L^{2}(\Gamma\backslash G)$ 的非零闭
$G$-不变子空间，则 $H$ 含有不可约子空间。}

\smallskip
\noindent\emph{证明（Langlands）.}\ $H$ 是 $G$-不变的，故每个 $T_{\phi}$ 在 $H$ 上诱导自同态。
先证：对适当选取的满足 (6) 的 $\phi$，$T_{\phi}$ 在 $H$ 上的限制非零。
若 $0\neq\xi\in H$，则当 $g$ 接近单位元时 $\rho(g)\xi$ 接近 $\xi$。
于是取 $\phi$ 满足 (6)、$\phi>0$、$\int_{G}\phi(g)dg=1$ 且支集不是一点，
则 $T_{\phi}\xi$ 接近 $\xi$，从而 $T_{\phi}$ 在 $H$ 上非零。

设 $L\subset H$ 是 $T_{\phi}$ 的非零特征值的特征空间。由定理 8，它有限维。
取 $L_{0}$ 为 $L$ 与 $H$ 的某个非零闭不变子空间之交，并使 $L_{0}$ 关于此性质极小。
设 $V$ 是 $H$ 的使 $L\cap V=L_{0}$ 的最小闭不变子空间。我们证明 $V$ 不可约。
否则设 $V_{1}$ 是 $V$ 的真非零闭不变子空间，$V_{2}$ 为其正交补，于是 $V=V_{1}\oplus V_{2}$。
取 $0\neq f\in L_{0}$，写 $f=f_{1}+f_{2}$，$f_{i}\in V_{i}$。由
$0=T_{\phi}f-\lambda f=(T_{\phi}f_{1}-\lambda f_{1})+(T_{\phi}f_{2}-\lambda f_{2})$
与 $T_{\phi}f_{i}-\lambda f_{i}\in V_{i}$，得 $T_{\phi}f_{i}-\lambda f_{i}=0$。
故 $f_{i}\in L\cap V_{i}$。由 $L_{0}$ 的极小性，$L_{0}=L\cap V_{i}$ 对某个 $i$ 成立，
譬如 $L_{0}=L\cap V_{1}$。这与 $V$ 的最小性矛盾。\hfill$\square$

\par\medskip\noindent\textbf{定理 11.}\quad\emph{$L^{2}(\Gamma\backslash G)$ 分解为闭不可约
子空间的直和，每个子空间给出 $G$ 的一个不可约容许表示。}

\smallskip
\noindent\emph{证明.}\ 由 Zorn 引理，取两两正交的闭不可约子空间之极大族 $S$。
令 $H=\bigoplus_{V\in S}V$。若 $H$ 是真子空间，则对它的正交补用引理 10，
与 $S$ 的极大性矛盾。容许性留给读者。\hfill$\square$

由定理 4，这些闭不可约子空间中每个至多含一个 $K$-不变向量。

定理 11 可以看作定理 9 的一个更令人满意的延伸。事实上，若 $\phi$ 是 $\Delta$ 的
一个出现在 $L^{2}(\Gamma\backslash\mathfrak{H})$ 中的特征函数，则它的右平移张成
$L^{2}(\Gamma\backslash G)$ 的一个不可约子空间。反之，$\Delta$ 在每个不可约子空间上
作用为纯量。若该子空间恰含一个 $K$-不变向量，这个向量就是定理 9 中基元素之一。

然而，$L^{2}(\Gamma\backslash G)$ 中确实有些不可约子空间不含 $K$-不变向量。
它们可以由\emph{全纯模形式}如下构造。设 $f:\Gamma\backslash\mathfrak{H}\to\mathbb{C}$
是全纯函数，满足
\[
f\!\left(\frac{az+b}{cz+d}\right)=(cz+d)^{k}f(z),\qquad
\begin{pmatrix}a&b\\c&d\end{pmatrix}\in\Gamma .
\]
若 $k$ 充分大，这样的 $f$ 总存在，这可从 Riemann--Roch 定理看出。
把 $f$ 看作 $G$ 上的函数时，它右 $K$-不变，但不左 $\Gamma$-不变，
故不是 $\Gamma\backslash G$ 上的函数。不过可以修改它，牺牲右 $K$-不变性以换取真正的左
$\Gamma$-不变性。定义
\[
F\begin{pmatrix}a&b\\c&d\end{pmatrix}=(ci+d)^{-k}f\!\left(\frac{ai+b}{ci+d}\right).
\]
则对 $\gamma\in\Gamma$ 有 $F(\gamma g)=F(g)$，而 $F(g\kappa_{\theta})=e^{ik\theta}F(g)$。
$F$ 张成的不可约表示不含 $K$-不变向量。这就是权为 $k$ 的\emph{全纯离散序列表示}。

表示论的进路还有另一种推广，即推广到阿代尔群上的自守形式。
设 $\Gamma=\mathrm{SL}(2,\mathbb{Z})$，$\A$ 是 $\mathbb{Q}$ 的阿代尔环。
嵌入 $\mathrm{SL}(2,\mathbb{R})\to\GL(2,\A)$ 在无穷位诱导同胚
\[
\Gamma\backslash G\longrightarrow \GL(2,\mathbb{Q})Z_{\A}\backslash\GL(2,\A)\Big/\prod_{p}K_{p},
\]
其中 $Z_{\A}$ 是 $\GL(2,\A)$ 的中心，$K_{p}=\GL(2,\mathbb{Z}_{p})$ 是 $\GL(2,\mathbb{Q}_{p})$
的极大紧子群。于是 $\Gamma\backslash G$ 上的函数可以重新解释为
$\GL(2,\mathbb{Q})Z_{\A}\backslash\GL(2,\A)$ 上的函数。特别地有连续嵌入
\[
L^{2}(\Gamma\backslash G)\hookrightarrow
L^{2}(\GL(2,\mathbb{Q})Z_{\A}\backslash\GL(2,\A)).
\]

现在，对 $L^{2}(\GL(2,\mathbb{Q})Z_{\A}\backslash\GL(2,\A))$ 的研究可以沿完全相同的路线进行。
只是积分算子类变大了：除环 $\mathcal{H}$ 之外，还有它的 $p$ 进类似物，即诸 \emph{Hecke 算子}环。
为看清这一点，固定素数 $p$。令 $\mathcal{H}_{p}$ 为 $\GL(2,\mathbb{Q}_{p})$ 上光滑（即局部常值）
紧支函数关于卷积构成的环。这个环非交换，但其中 $K_{p}$-双不变函数构成的子环
$\mathcal{H}_{p}^{\circ}$ 交换，可与定理 3 对照。例如，$K_{p}\binom{p\ 0}{0\ 1}K_{p}$
的特征函数就是此环的一个元素，它是经典 Hecke 算子 $T_{p}$ 的阿代尔化，
而 $T_{p}$ 从 Hecke 本征形式中提取第 $p$ 个 Fourier 系数。
最后一节我们研究这类算子时会回到这个观点。

\begin{annote}
定理 3 在 $p$ 进情形即 Satake 同构：$\mathcal{H}_{p}^{\circ}$ 同构于
$\mathbb{C}[X_{*}(T)]^{W}$，即余特征标格在 Weyl 群下的不变量环；对 $\GL(2)$ 它就是
$\mathbb{C}[t+t^{-1}]$。因此局部球表示由半单共轭类即 Satake 参数分类，
而整体 Hecke 特征值经归一化后正是 $L$ 函数的局部因子。原书把 $\mathcal{H}_{p}^{\circ}$
的交换性归到 $\GL(2)$ 的直接计算，一般约化群的情形要等到 Satake 1963 年的定理。
\par\vspace{0.15em}{\raggedleft 【来源：Bump, \emph{Automorphic Forms and Representations}, Cambridge, 1997, Ch.~4；Satake, \emph{Theory of spherical functions on reductive algebraic groups over $p$-adic fields}, Publ. Math. IH\'ES 18 (1963)（A）】\par}
\end{annote}

\section{Green 函数与谱预解式}

本节的参考文献是 Hejhal [21, Section 6] 与 Bump [6, Chapter 2, Section 3]。
我们关心 $\mathfrak{H}$ 上满足下述条件的函数 $f$。

\par\medskip\noindent\textbf{性质 S.}\quad\emph{函数 $f$ 左 $K=\mathrm{SO}(2)$-不变，
在 $K$ 的不动点 $i$ 处可以奇异，并且是 Laplacian 的特征函数。}

我们将看到，对 $\Delta$ 的每个特征值 $\lambda$，有两个线性无关的此类函数：
一个（\emph{Green 函数}）在边界处性状良好，另一个（\emph{球函数}）在 $i$ 处连续。
两者各有用途。由于 $\mathfrak{H}\cong G/K$，它们可以看作 $K\backslash G/K$ 上的函数
（在 $K$ 上可能无定义），并且是 Laplace--Beltrami 算子的特征函数。

用 Cayley 变换 $z\mapsto w=(z-i)/(z+i)$ 把上半平面映到单位圆盘 $\mathfrak{D}$。
令 $r=|w|$。性质 S 中的 $f$ 左 $K$-不变，故 $f(z)$ 只依赖 $r$。
记 $f(z)=W(r)$。于是满足性质 S 的函数 $f$ 由 $(0,1)$ 上的函数 $W$ 决定，二者关系为
\[
W(r)=f\begin{pmatrix}y^{1/2}&\\&y^{-1/2}\end{pmatrix},
\]
其中 $r=(y-1)/(y+1)\in(0,1)$。特征函数这一要求等价于微分方程
\begin{equation}\tag{7}
W''(r)+\frac{1}{r}W'(r)+\frac{4\lambda}{(1-r^{2})^{2}}W(r)=0 .
\end{equation}

这个微分方程在 $0$ 与 $1$ 处有正则奇点，有两个值得注意的解：
一个在 $0$ 处性状良好，另一个在 $1$ 处性状良好。
（我们假定读者熟悉二阶线性微分方程的正则奇点，特别是指标方程，
见 Whittaker 与 Watson [52, Section 10.3]。）本节关心在边界、即 $r=1$ 处性状良好的解，
另一个解留到下一节。

设 $\lambda=-s$，其中 $s$ 是正实数。在 $r=1$ 处，(7) 的指标方程的根是
$\tfrac12(1\pm\sqrt{1+4s})$。只有一个是正的，故在边界附近为零的 (7) 的解 $g$ 除常数倍外唯一。
可以用它来研究 Laplacian 的预解式。

\par\medskip\noindent\textbf{引理 12.}\quad\emph{$g_{s}(r)$ 在 $r=0$ 处有对数奇点。}

\smallskip
\noindent\emph{证明.}\ $r=0$ 处指标方程的根是 $0$，重数为 $2$，故一个解有对数奇点，
另一个解解析。若 $g_{s}$ 没有对数奇点，则 $g_{s}(r)$ 在 $[-1,1]$ 上实解析，因而有最大值或最小值。
在这样的点上 $g_{s}'(r)=0$；又 $\lambda=-s<0$，由 (7) 知 $g_{s}$ 与 $g_{s}''$ 同号，
这在最大值或最小值处不可能，因为 $g_{s}(-1)=g_{s}(1)=0$。\hfill$\square$

置
\[
g_{s}(z,\zeta)=g_{s}\!\left(\left|\frac{z-\zeta}{z-\bar\zeta}\right|\right),
\qquad z,\zeta\in\mathfrak{H}.
\]
这就是一个 \emph{Green 函数}。

\par\medskip\noindent\textbf{定理 13.}\quad
\begin{equation}\tag{8}
\begin{aligned}
&\left[-y^{2}\left(\frac{\partial^{2}}{\partial x^{2}}+\frac{\partial^{2}}{\partial y^{2}}\right)+s\right]g(z,\zeta)=0;\\
&g_{s}(z,\zeta)\ \text{在角线}\ z=\zeta\ \text{上奇异};\\
&g_{s}(z,\zeta)\to0\quad(y\to0);\\
&g_{s}(z,\zeta)=g_{s}(\zeta,z);\\
&g_{s}(h(z),h(\zeta))=g_{s}(z,\zeta),\qquad h\in\mathrm{SL}(2,\mathbb{R}).
\end{aligned}
\end{equation}

\smallskip
更完整的细节见 Bump [6, Proposition 2.3.4, p.\ 181]。这里没有什么是特别意外的：
例如第一条性质归结为 (7)，第二条来自引理 12，第三条来自 $g_{s}$ 的边界性状，
最后两条则因为 $|(z-\zeta)/(z-\bar\zeta)|$ 在 $z$ 与 $\zeta$ 互换、或对二者同时作用
$\mathrm{SL}(2,\mathbb{R})$ 的元素时不变。\hfill$\square$

$g_{s}$ 在 $0$ 处有对数奇点，故可归一化使 $g_{s}(r)-\frac{1}{2\pi}\log r$ 当 $r\to0$ 时有界。
于是 $g_{s}'(r)-\frac{1}{2\pi r}$ 在 $r=0$ 附近解析。

\par\medskip\noindent\textbf{定理 14.}\quad\emph{若 $f\in C_{c}^{\infty}(\mathfrak{H})$，则
（写 $\zeta=\xi+i\eta$）}
\[
\int_{\mathfrak{H}}g_{s}(z,\zeta)\left[-\eta^{2}\left(\frac{\partial^{2}}{\partial\xi^{2}}
+\frac{\partial^{2}}{\partial\eta^{2}}\right)+s\right]f(\zeta)\,\frac{d\xi\wedge d\eta}{\eta^{2}}=f(z).
\]

\smallskip
\noindent\emph{证明.}\ 更完整的细节见 Bump [6, Proposition 2.3.4, p.\ 181]。这里简述其证明。
令 $w=(z-\zeta)/(z-\bar\zeta)=u+iv\in\mathfrak{D}$，并令 $F:\mathfrak{D}\to\mathbb{C}$ 由
$F(w)=f(z)$ 定义。在 $w$ 坐标下，我们要证
\[
\int_{\mathfrak{D}}g_{s}(|w|)\left[-\left(\frac{\partial^{2}}{\partial u^{2}}
+\frac{\partial^{2}}{\partial v^{2}}\right)+\frac{4s}{(1-|w|^{2})^{2}}\right]F(w)\,du\wedge dv=F(0).
\]
设 $B_{r}$ 是半径为 $r$ 的圆盘，$R<1$ 充分大，使 $F$ 的支集含于 $B_{R}$。左端等于
\[
\lim_{\epsilon\to0}\int_{B_{R}-B_{\epsilon}}g_{s}(|w|)
\left[-\left(\frac{\partial^{2}}{\partial u^{2}}+\frac{\partial^{2}}{\partial v^{2}}\right)
+\frac{4s}{(1-|w|^{2})^{2}}\right]F(w)\,du\wedge dv .
\]
像命题 2 那样用 Stokes 定理，它等于
\[
\begin{aligned}
&\lim_{\epsilon\to0}\int_{B_{R}-B_{\epsilon}}F(w)\left[-\left(\frac{\partial^{2}}{\partial u^{2}}
+\frac{\partial^{2}}{\partial v^{2}}\right)+\frac{4s}{(1-|w|^{2})^{2}}\right]g_{s}(|w|)\,du\wedge dv\\
&\quad+\lim_{\epsilon\to0}\int_{C_{\epsilon}}F(w)\left(\frac{\partial g_{s}(|w|)}{\partial u}\,dv
-\frac{\partial g_{s}(|w|)}{\partial v}\,du\right)\\
&\quad-\lim_{\epsilon\to0}\int_{C_{\epsilon}}g_{s}(|w|)\left(\frac{\partial F(w)}{\partial u}\,dv
-\frac{\partial F(w)}{\partial v}\,du\right),
\end{aligned}
\]
其中 $C_{\epsilon}$ 是绕原点逆时针、半径为 $\epsilon$ 的圆周。
（本来还有沿 $C_{R}$ 的积分项，但 $C_{R}$ 落在 $F$ 的支集之外，故为零。）

由定理 13，第一项为零。最后一项为零，因为弧长收缩比 $g_{s}$ 的对数发散更快。
为计算中间项，令 $w=re^{i\theta}$。由链式法则
\[
\frac{\partial g_{s}(|w|)}{\partial u}\,dv-\frac{\partial g_{s}(|w|)}{\partial v}\,du
=r\,g_{s}'(r)\,d\theta .
\]
于是得到
\[
\lim_{\epsilon\to0}\int_{0}^{2\pi}F(\epsilon e^{i\theta})\,d\theta\,\epsilon g_{s}'(\epsilon)=F(0),
\]
因为当 $\epsilon\to0$ 时 $g_{s}'(\epsilon)\sim1/(2\pi\epsilon)$。\hfill$\square$

\begin{annote}
令 $\lambda=\tfrac14+t^{2}$ 并置 $x=(1+r^{2})/(1-r^{2})$，则 (7) 的两个解分别是第一类与
第二类 Legendre 函数 $P_{-1/2+it}(x)$ 与 $Q_{-1/2+it}(x)$。在 $i$ 处正则的那个（第 4 节的球函数）
对应 $P$，在边界 $r=1$ 处消失的那个（Green 函数 $g_{s}$）对应 $Q$。
引理 12 说的正是 $Q$ 在 $r=0$ 处有对数奇点而 $P$ 解析，这也就解释了为何 Green 函数
不能取成球函数。
\par\vspace{0.15em}{\raggedleft 【来源：Hejhal, Vol.~2, Ch.~1；Whittaker--Watson, \emph{A Course of Modern Analysis}, 4th ed., Cambridge, 1927, \S14.5（A）】\par}
\end{annote}


\par\medskip\noindent\textbf{命题 15.}\quad\emph{级数}
\[
G_{s}(z,\zeta)=\sum_{\gamma\in\{\pm1\}\backslash\Gamma}g_{s}(z,\gamma(\zeta))
=\sum_{\gamma\in\{\pm1\}\backslash\Gamma}g_{s}(\gamma(z),\zeta)
\]
\emph{在 $z$ 与 $\zeta$ 不 $\Gamma$-等价时绝对收敛。}

\smallskip
见 Bump [6, Proposition 2.3.5, p.\ 184]。\hfill$\square$

$G_{s}(z,\zeta)$ 是\emph{自守 Green 函数}。我们将看到，它是 Laplacian 的预解式的积分核。

\par\medskip\noindent\textbf{定理 16.}\quad\emph{$G_{s}$ 对一切 $(z,\zeta)$ 有定义且实解析，
例外是对某个 $\gamma\in\Gamma$ 有 $\zeta=\gamma(z)$ 的情形。我们有}
\[
\begin{gathered}
\left[-y^{2}\left(\frac{\partial^{2}}{\partial x^{2}}+\frac{\partial^{2}}{\partial y^{2}}\right)+s\right]G_{s}(z,\zeta)=0;\\
G_{s}(h(z),h(\zeta))=G_{s}(z,\zeta),\qquad h\in G,\\
G_{s}(z,\zeta)=G_{s}(\gamma(z),\zeta)=G_{s}(z,\gamma(\zeta)),\qquad\gamma\in\Gamma,\\
G_{s}(z,\zeta)=\overline{G_{s}(\zeta,z)},
\end{gathered}
\]
\emph{并且在 $z=\zeta$ 附近有 $G_{s}(z,\zeta)\sim\frac{e}{2\pi}\log|z-\zeta|$，
其中 $e$ 是 $\zeta$ 在 $\Gamma$ 中的稳定化子群的阶。又对 $f\in C^{\infty}(\Gamma\backslash\mathfrak{H})$，}
\begin{equation}\tag{9}
\int_{\Gamma\backslash\mathfrak{H}}G_{s}(z,\zeta)\left[-\eta^{2}
\left(\frac{\partial^{2}}{\partial\xi^{2}}+\frac{\partial^{2}}{\partial\eta^{2}}\right)+s\right]
f(\zeta)\,\frac{d\xi\wedge d\eta}{\eta^{2}}=f(z)\qquad(\zeta=\xi+i\eta).
\end{equation}

\smallskip
见 Bump [6, Proposition 2.3.5] 与 Hejhal [21, Proposition 6.5, p.\ 33]。

\smallskip
\noindent\emph{证明.}\ 这些性质大多由 $g_{s}$ 的相应性质推出。我们只证 (9)。
需要一个函数 $F\in C_{c}^{\infty}(\mathfrak{H})$，使 $f(z)=\sum_{\gamma\in\Gamma}F(\gamma z)$。
为构造 $F$，取 $\mathfrak{H}$ 上光滑、非负、支集紧且包含 $\Gamma$ 的基本区域的函数 $u$。
则对一切 $z$，函数 $\sum_{\gamma\in\Gamma}u(\gamma z)$ 为正；当 $z$ 限制在紧集上时此和有限。
于是它是光滑的正值函数，可以用它作除数。现在
\[
F(z)=\frac{u(z)f(z)}{\sum_{\gamma\in\Gamma}u(\gamma z)}
\]
就具有所需性质。

把它与 $G_{s}$ 的定义代入，并用 (8)，得到
\[
\sum_{\gamma,\delta\in\Gamma}\int_{\Gamma\backslash\mathfrak{H}}g_{s}(\delta(z),\gamma(\zeta))
\left[-\eta^{2}\left(\frac{\partial^{2}}{\partial\xi^{2}}+\frac{\partial^{2}}{\partial\eta^{2}}\right)+s\right]
F(\gamma(\zeta))\,\frac{d\xi\wedge d\eta}{\eta^{2}} .
\]
其中一个求和可以与积分合并，给出
\[
\begin{aligned}
\sum_{\gamma\in\Gamma}\int_{\mathfrak{H}}g_{s}(\delta(z),\zeta)
\left[-\eta^{2}\left(\frac{\partial^{2}}{\partial\xi^{2}}+\frac{\partial^{2}}{\partial\eta^{2}}\right)+s\right]
F(\zeta)\,\frac{d\xi\wedge d\eta}{\eta^{2}}
&=\sum_{\delta}F(\delta(z))\\
&=f(z).
\end{aligned}
\]
这就完成了证明。\hfill$\square$

设 $s>0$ 为常数。容易验证，角线上的对数奇点不足以使积分发散，即
\[
\int_{\Gamma\backslash\mathfrak{H}}\int_{\Gamma\backslash\mathfrak{H}}|G_{s}(z,\zeta)|^{2}
\,\frac{dx\wedge dy}{y^{2}}\,\frac{d\xi\wedge d\eta}{\eta^{2}}<\infty .
\]
于是相应的积分算子是 Hilbert--Schmidt 的，我们把它记作 $R(s,\Delta)$。
它是 Laplacian 的\emph{预解式}。若 $\phi$ 是 $\Delta$ 的特征函数、特征值为 $\lambda$，
则由 (9) 知它也是 $R(s,\Delta)$ 的特征函数，特征值为 $(\lambda+s)^{-1}$。

\par\medskip\noindent\textbf{定理 17.}\quad
(i) \emph{$\Delta$ 在 $L^{2}(\Gamma\backslash\mathfrak{H})$ 上的特征值 $\lambda_{i}$ 趋于无穷，
且满足 $\sum\lambda_{i}^{-2}<\infty$。（求和时略去常数函数对应的特征值 $\lambda_{0}=0$。）}

(ii) \emph{Laplacian $\Delta$ 可延拓为 Hilbert 空间 $L^{2}(\Gamma\backslash\mathfrak{H})$ 上的自伴算子。}

\smallskip
\noindent\emph{证明.}\ 设 $\operatorname{Re}(s)>0$。$R(s,\Delta)$ 是 Hilbert--Schmidt 的，
故 $\sum_{i}(\lambda_{i}+s)^{-2}<\infty$，从而 $\sum_{i}\lambda_{i}^{-2}<\infty$。
又由定理 9，取 $\phi_{i}$ 为 $H=L^{2}(\Gamma\backslash\mathfrak{H})$ 的由 $\Delta$ 的特征向量
构成的基，相应特征值为 $\lambda_{i}$。

下证 (ii)。令 $\mathfrak{D}_{\Delta}$ 为 $\mathfrak{H}$ 上形如 $\sum a_{i}\phi_{i}$、
且满足 $\sum\lambda_{i}^{2}|a_{i}|^{2}<\infty$ 的元素构成的线性子空间。在此空间上定义
\[
\Delta\Bigl(\sum a_{i}\phi_{i}\Bigr)=\sum\lambda_{i}a_{i}\phi_{i}.
\]
由于 $\lambda_{i}$ 趋于无穷，特别地与零有正距离，不难验证这个算子闭，且事实上自伴。
证毕。\hfill$\square$

我们把 $\Delta$ 的特征值按递增次序排列：
\[
\lambda_{0}=0<\lambda_{1}\le\lambda_{2}\le\lambda_{3}\le\cdots .
\]
特征值 $\lambda_{0}$ 对应常数函数，重数恰为 $1$。落在 $(0,\tfrac14)$ 中的特征值称为
\emph{例外特征值}，它们的性质有本质的不同。例如，对应例外特征值的球函数不是缓增的：
它们比对应 $\lambda\ge\tfrac14$ 的球函数增长更快（见第 4 节）。
例外特征值对应 Selberg zeta 函数在实轴上 $0$ 与 $1$ 之间的零点。
与此相对，对应 $\lambda\ge\tfrac14$ 的零点满足 Riemann 假设（见第 7 节）。
Randol [39] 证明了：对某些 $X$，例外特征值确实出现。另一方面，Selberg [44] 猜想：
$\mathrm{SL}(2,\mathbb{Z})$ 的同余子群的尖点谱中没有例外特征值；
由此还可推出，与四元数除代数相联的紧商 $\Gamma\backslash\mathfrak{H}$ 的谱中也没有例外特征值。

\begin{annote}
例外特征值等价于 Selberg zeta 函数 $Z(s)$ 在 $(1/2,1)$ 内有零点，见第 7 节定理 35。
Selberg 猜想断言 $\mathrm{SL}(2,\mathbb{Z})$ 的同余子群没有例外特征值，即 $\lambda_{1}\ge1/4$，
至今未决；目前最好的下界 $\lambda_{1}\ge975/4096$ 由 Kim 与 Sarnak 得到。
Randol 的例子说明紧曲面上小特征值确实可能出现，故猜想只对算术情形成立。
\par\vspace{0.15em}{\raggedleft 【来源：Selberg, J. Indian Math. Soc. 20 (1956)；Randol, Bull. Amer. Math. Soc. 80 (1974)；Kim--Sarnak, J. Amer. Math. Soc. 16 (2003)（A）】\par}
\end{annote}

\section{球函数}

如前所述，性质 S 归结为一个二阶微分方程，它有两个独立解。
一个在边界处性状良好，即 Green 函数；另一个在 $i$ 处性状良好，即所谓球函数。
引理 12 的实质就是：这两个解并不相同。本节研究球函数。

\par\medskip\noindent\textbf{定义.}\quad 设 $\lambda$ 是复常数。$\mathrm{SL}(2,\mathbb{R})$ 上的函数
$\sigma$ 若光滑、$K$-双不变，且 $\Delta\sigma=\lambda\sigma$，则称它为 \emph{$\lambda$-球函数}。

这里 $\Delta$ 是 Laplace--Beltrami 算子。在证明这样的函数存在且除常数倍外唯一之前，
我们先简要说明它如何嵌入表示论。

设 $(\pi,V)$ 是 $G$ 的不可约表示，$(\hat\pi,\hat V)$ 是其反步表示。
于是存在不变双线性配对 $\langle\,,\,\rangle:V\times\hat V\to\mathbb{C}$。
由定理 4，若 $V^{K}$ 与 $\hat V^{K}$ 非零，则它们都是一维的，我们假定如此。
（事实上只要一个非零，另一个也非零。）取非零的 $K$-不变向量
$v^{\circ}\in V^{K}$ 与 $\hat v^{\circ}\in\hat V^{K}$。函数
\[
\sigma(g)=\langle\pi(g)v^{\circ},\hat v^{\circ}\rangle
\]
显然 $K$-双不变。此外，把 $\Delta$ 看作泛包络代数 $U(\mathfrak{g})$ 的中心元素，
它在 $V$ 上作用为纯量 $\lambda$，而 $\sigma$ 继承了这一性质，故它是球函数。

下面不再借助这一构造，直接证明球函数存在且唯一。

\par\medskip\noindent\textbf{定理 18.}\quad
(i) \emph{设 $\lambda\in\mathbb{C}$。则存在唯一的 $\mathrm{SL}(2,\mathbb{R})$ 上光滑
$K$-双不变函数 $\omega_{\lambda}$，使 $\Delta\omega_{\lambda}=\lambda\omega_{\lambda}$，
$\omega_{\lambda}(1)=1$。}

(ii) \emph{若 $f:G\to\mathbb{C}$ 是任一满足 $\Delta f=\lambda f$ 的光滑函数，则}
\begin{equation}\tag{10}
\int_{K\times K}f(gkg')\,dk\,dk'=f(1)\,\omega_{\lambda}(g).
\end{equation}

(iii) \emph{若 $f$ 右 $K$-不变且 $\Delta f=\lambda f$，则}
\begin{equation}\tag{11}
\int_{K}f(hkg)\,dk=f(h)\,\omega_{\lambda}(g).
\end{equation}

\smallskip
\noindent\emph{证明.}\ 要使 $\Delta\omega=\lambda\omega$，需要
\[
W(r)=\omega\begin{pmatrix}y^{1/2}&\\&y^{-1/2}\end{pmatrix},
\qquad r=\frac{y-1}{y+1}
\]
满足 (7)。如前所见，这个微分方程在原点有正则奇点，一个解有界，另一个有对数奇点。
故 $\omega_{\lambda}$ 若存在则唯一。

为证存在性，取 $G$ 上任一连续函数 $f$，它是 $\Delta$ 的特征函数。
(10) 的左端是 $K$-双不变函数，也是 $\Delta$ 的特征函数。若 $f(1)=1$，它就满足 (i)，
存在性得证。例如可取 $f=f_{s}$，其中 $\lambda=s(1-s)$，且
\begin{equation}\tag{12}
f_{s}\!\left(\begin{pmatrix}y^{1/2}&*\\&y^{-1/2}\end{pmatrix}k\right)=y^{s},
\qquad y>0,\ k\in K .
\end{equation}

现在 $\omega$ 的存在性与唯一性都已确立。注意 (10) 与 (11) 的左端都是光滑的
$K$-双不变 $\Delta$-特征函数，故都是 $\omega_{\lambda}$ 的常数倍。
两种情形下都令 $g=1$ 即可定出常数。这就证明了 (ii) 与 (iii)。\hfill$\square$

\par\medskip\noindent\textbf{定理 19.}\quad\emph{设 $f$ 是 $G$ 上光滑、右 $K$-不变的函数，
且 $\Delta f=\lambda f$。则对 $\phi\in\mathcal{H}^{\circ}$ 有 $T_{\phi}f=\chi_{\lambda}(\phi)f$，其中}
\begin{equation}\tag{13}
\chi_{\lambda}(\phi)=\int_{G}\phi(g)\,\omega_{\lambda}(g)\,dg .
\end{equation}

\smallskip
\noindent\emph{证明.}\ 由定理 18，
\begin{equation}\tag{14}
\int_{K}f(hkg)\,dk=f(h)\,\omega_{\lambda}(g).
\end{equation}
注意 $T_{\phi}f$ 是 $f$ 的右平移的平均，而右平移与左平移交换。
于是可对 (14) 两端作用 $T_{\phi}$，得到
\[
\int_{K}(T_{\phi}f)(hkg)\,dk=f(h)\,(T_{\phi}\omega_{\lambda})(g).
\]
在此恒等式中取 $g=1$。$T_{\phi}f$ 右 $K$-不变，故 $g=1$ 时左端的被积函数为常数，
左端就等于 $(T_{\phi}f)(h)$。另一方面，$(T_{\phi}\omega_{\lambda})(1)$ 等于积分 (13)，
故 $T_{\phi}f(h)=\chi_{\lambda}(\phi)f(h)$。\hfill$\square$

\par\medskip\noindent\textbf{定理 20.}\quad\emph{若 $\phi_{1},\phi_{2}\in\mathcal{H}^{\circ}$，则}
\[
\chi_{\lambda}(\phi_{1}*\phi_{2})=\chi_{\lambda}(\phi_{1})\,\chi_{\lambda}(\phi_{2}).
\]

\smallskip
\noindent\emph{证明.}\ 任取特征函数 $f$（例如 (12) 中的 $f_{s}$），对 $T_{\phi}$ 用定理 19
即可。\hfill$\square$

于是 $\chi_{\lambda}:\mathcal{H}^{\circ}\to\mathbb{C}$ 是 $\mathcal{H}^{\circ}$ 的\emph{特征}。
我们回到本节开头引入的表示论观点，以解释 $\chi(\phi)$ 的含义。
先回顾 $\mathrm{SL}(2,\mathbb{R})$ 的一类不可约表示的构造与参数化，即\emph{主系列}。

设 $s$ 是复常数。以 $P_{s}^{+}$ 记满足下式的函数 $f:G\to\mathbb{C}$ 之集：
\begin{equation}\tag{15}
f\!\left(\begin{pmatrix}y^{1/2}&xy^{-1/2}\\&y^{-1/2}\end{pmatrix}g\right)
=(-1)^{\epsilon}\operatorname{sgn}(y)\,y^{s}f(g),
\end{equation}
其中 $\epsilon=0$，并要求 $f$ 在 $K$ 上的限制平方可积。类似地，以 $P_{s}^{-}$ 记满足 (15)、
$\epsilon=1$ 的函数空间。$G$ 通过右平移作用在这些空间上，记相应表示为
$(\pi_{s}^{\pm},P_{s}^{\pm})$。除 $2s$ 为整数、且对 $\pi_{s}^{+}$ 为偶数、对 $\pi_{s}^{-}$ 为奇数
的情形外，这些表示不可约。我们把 $\pi_{s}^{+}$ 与 $\pi_{s}^{-}$ 分别称为
\emph{奇主系列}与\emph{偶主系列}表示。

\begin{annote}
此处命名与 (15) 的取法及第 2 节的约定都不一致：$P_{s}^{+}$ 取 $\epsilon=0$，
而由 (15) 右端的 $\operatorname{sgn}(y)^{\epsilon}$ 可知，决定奇偶的是中心元 $-I$ 的特征值
$(-1)^{\epsilon}$，故 $\pi_{s}^{+}$ 应为\emph{偶}主系列。原书同页后文又说“只有偶主系列
$\pi_{s}^{+}$ 含 $K$-不变向量”，也把 $\pi_{s}^{+}$ 当作偶的；可见“奇、偶”二字在此处
被写颠倒了。译文照录原文，读者按 $\pi_{s}^{+}$ 为偶系列理解，并注意 $k\equiv\epsilon\bmod2$、
$\pm=(-1)^{\epsilon}$。
\par\vspace{0.15em}{\raggedleft 【来源：Bump, \emph{Automorphic Forms and Representations}, 1997, \S2.6 与 \S3.3；原书 p.~168 前后两处自相矛盾（A）】\par}
\end{annote}


设 $\operatorname{re}(s)>\tfrac12$。缠绕积分 $M(s):P_{s}^{\pm}\to P_{1-s}^{\pm}$ 由
\begin{equation}\tag{16}
(M(s)f)(g)=\int_{-\infty}^{\infty}
f\!\left(\begin{pmatrix}&-1\\1&\end{pmatrix}
\begin{pmatrix}1&x\\&1\end{pmatrix}g\right)dx
\end{equation}
给出。可以验证此积分绝对收敛。$M(s)$ 经解析延拓成为对一切使 $2s$ 不为与 $\epsilon$
模 $2$ 同余的整数的 $s$ 都有定义的缠绕映射。其实我们只断言它延拓为无穷小等价，
即底层的 $(\mathfrak{g},K)$-模的同构；但在重要的特殊情形 $\operatorname{re}(s)=\tfrac12$，
它是真正的同构。于是 $\pi_{1/2+it}^{\pm}\cong\pi_{1/2-it}^{\pm}$。

设 $k\equiv\epsilon\ \mathrm{mod}\ 2$。置
\begin{equation}\tag{17}
f_{s,k}\!\left(\begin{pmatrix}y^{1/2}&x\\&y^{-1/2}\end{pmatrix}\kappa_{\theta}\right)
=\operatorname{sgn}(y)^{\epsilon}\,y^{s}e^{ik\theta}.
\end{equation}
这是 $P_{s}^{\pm}$ 中的一个 $K$-有限向量，其中 $\pm$ 取 $(-1)^{\epsilon}$。
令 $\tilde f_{1-s,k}=M(s)f_{s,k}$。直接计算可得
\begin{equation}\tag{18}
\tilde f_{k,1-s}=(-i)^{k}\sqrt{\pi}\,
\frac{\Gamma(s)\Gamma\left(s-\frac12\right)}
{\Gamma\left(s+\frac{k}{2}\right)\Gamma\left(s-\frac{k}{2}\right)}\,f_{k,1-s}.
\end{equation}
见 Bump [6, Proposition 2.6.3]。

$t$ 为实数时，表示 $\pi_{1/2+it}^{\pm}$ 是酉的，称为\emph{酉主系列}表示。
此外，$s$ 为 $0$ 与 $1$ 之间的实数时 $\pi_{s}^{+}$ 也是酉的，这些是\emph{补系列}表示。
它们对应自守谱中 Laplacian 的例外特征值。此外还有少数其他不可约表示，
即离散序列（与全纯模形式有关）与平凡表示。但只有偶主系列 $\pi_{s}^{+}$ 含 $K$-不变向量。
这些事实的证明都在 Bump [6, Chapter 2]。

$P_{s}^{+}$ 中的 $K$-不变向量正是 (12) 中的函数 $f_{s}$。
现在若 $\phi\in\mathcal{H}^{\circ}$，则因 $\phi$ 是 $K$-双不变的，算子 $\pi^{+}(\phi)$
把 $P_{s}^{+}$ 映到其唯一的 $K$-不变向量 $f_{s}$ 上。
于是秩一算子 $\pi_{s}^{+}(\phi)$ 的迹就是它在 $f_{s}$ 上的特征值，故
\begin{equation}\tag{19}
\operatorname{tr}\pi_{s}^{+}(\phi)=\chi_{\lambda}(\phi).
\end{equation}
\section{Plancherel 公式}

设 $\phi\in\mathcal{H}^{\circ}$。定义
\begin{equation}\tag{20}
g(u)=e^{u/2}\int_{-\infty}^{\infty}
\phi\!\left(\begin{pmatrix}e^{u/2}&\\&e^{-u/2}\end{pmatrix}
\begin{pmatrix}1&x\\&1\end{pmatrix}\right)dx,
\end{equation}
并令
\begin{equation}\tag{21}
h(t)=\int g(u)\,e^{iut}\,dt
\end{equation}
为它的 Fourier 变换。

\begin{annote}
校勘：(21) 原书印作 $h(t)=\int g(u)e^{iut}dt$，积分变量应为 $du$；
同页证明中又印作 $\int_{\infty}^{\infty}g(u)e^{itu}du$，上下限均应为一 $\infty$ 与 $\infty$。
译文照录原文，读者按 $du$ 理解即可。
\par\vspace{0.15em}{\raggedleft 【来源：原书 p.~169，式 (21) 及其证明（A）】\par}
\end{annote}


\par\medskip\noindent\textbf{定理 21.}\quad\emph{函数 $g$ 与 $h$ 都是偶函数，且 $g$ 紧支。
若 $\lambda=\tfrac14+t^{2}$，则}
\begin{equation}\tag{22}
\chi_{\lambda}(\phi)=h(t).
\end{equation}

\smallskip
\noindent\emph{证明.}\ 取 (12) 中的 $f_{s}$，令 $s=\tfrac12+it$，于是
$\lambda=\tfrac14+t^{2}=s(1-s)$。由定理 18，
\[
\int_{K}f_{s}(kg)\,dk=\omega_{\lambda}(g),
\]
又
\[
\chi_{\lambda}(\phi)=\int_{G}\int_{K}\phi(g)f_{s}(gk)\,dk\,dg .
\]
交换积分次序，并作变量替换 $g\to gk^{-1}$（注意 $\phi$ 右 $K$-不变），得到
\[
\chi_{\lambda}(\phi)=\int_{G}\phi(g)f_{s}(g)\,dg .
\]
现在采用坐标
\begin{equation}\tag{23}
g=\begin{pmatrix}e^{u/2}&\\&e^{-u/2}\end{pmatrix}
\begin{pmatrix}1&x\\&1\end{pmatrix}\kappa_{\theta},
\qquad dg=\frac{1}{2\pi}\,du\,dx\,d\theta .
\end{equation}
注意在这些坐标下 $f_{s}(g)=e^{u/2}e^{iut}$，于是得到
\[
\chi_{\lambda}(\phi)=\int_{-\infty}^{\infty}g(u)\,e^{itu}\,du,
\]
这就证明了 (22)。我们指出，若 $s\to1-s$、即 $t\to-t$，则 $\omega_{\lambda}$ 不变，
从而特征 $\chi_{\lambda}(\phi)$ 不变。于是 (22) 蕴涵 $h$ 是偶函数，由 Fourier 反演，
$g$ 也是偶函数。\hfill$\square$

\par\medskip\noindent\textbf{定理 22.}\quad
(i) \emph{有}
\[
\phi(1)=\frac{1}{2\pi}\int_{0}^{\infty}t\,h(t)\tanh(\pi t)\,dt .
\]
(ii) \emph{若 $\phi_{1},\phi_{2}$ 是 $K$-不变且紧支的，则}
\[
\int_{G}\phi_{1}(g)\overline{\phi_{2}(g)}\,dg
=\frac{1}{2\pi}\int_{0}^{\infty}t\,h_{\phi_{1}}(t)\overline{h_{\phi_{2}}(t)}
\tanh(\pi t)\,dt .
\]

\smallskip
证明见 Knapp [27, Chapter 11]、Gelfand--Graev--Piatetski-Shapiro [16, Chapter 2, Section 6]
与 Varadarajan [47, Theorem 39, p.\ 205]。下面我们在若干引理之后给出一个不用 Lie 理论的证明。
论证的一部分与 Hejhal [20, Proposition 4.1, p.\ 15] 平行，记号也与他一致。

定理 22 的 (i) 是非交换群 $\mathrm{SL}(2,\mathbb{R})$ 上的 Fourier 反演公式。
它有时被称为 \emph{Plancherel 公式}，因为它蕴涵 (ii)，而 (ii) 才是真正的 Plancherel 公式。
测度 $\frac{1}{2\pi}t\tanh(\pi t)\,dt$ 称为偶酉主系列上的 \emph{Plancherel 测度}。
（既然我们只考虑 $\phi\in\mathcal{H}^{\circ}$，就不需要其他不可约酉表示。）
Plancherel 测度与缠绕积分 (16) 密切相关。事实上，
$M(\tfrac12-it)\circ M(\tfrac12+it)$ 是不可约空间 $P_{1/2+it}^{+}$ 的 $G$-等变自同态，
故由 Schur 引理是纯量。用 (18) 可以验证，这个纯量的倒数是 $\frac{1}{\pi}t\tanh(\pi t)$。
见 Knapp 与 Stein [28]。

\par\medskip\noindent\textbf{引理 23.}\quad\emph{设 $\eta\ge y>1$。则}
\[
\phi\!\left(\begin{pmatrix}1&x\\&1\end{pmatrix}
\begin{pmatrix}y^{1/2}&\\&y^{-1/2}\end{pmatrix}\right)
=\phi\begin{pmatrix}\eta^{1/2}&\\&\eta^{-1/2}\end{pmatrix},
\]
\emph{其中}
\[
x=\pm\sqrt{y(\eta+\eta^{-1}-y-y^{-1})}.
\]

\smallskip
\noindent\emph{证明.}\ 把 $K\backslash G/K$ 与 $K\backslash\mathfrak{H}$ 等同，我们要找 $x$，
使 $\eta i$ 与 $z=x+iy$ 这两个矩阵在 $\mathfrak{H}$ 中的像属于同一 $SO(2)$ 轨道。
用 Cayley 变换 $z\to(z-i)/(z+i)$ 把 $\mathfrak{H}$ 映到单位圆盘，两者的像到原点的距离必须相等。即
\[
\frac{\eta-1}{\eta+1}=\left|\frac{z-i}{z+i}\right|=\sqrt{\frac{a-2y}{a+2y}},
\qquad a=x^{2}+y^{2}+1 .
\]
对此式两端作用映射 $t\to(1+t^{2})/(1-t^{2})$，得 $\eta+\eta^{-1}=a/y$。
解此方程即得 $x$。\hfill$\square$

在 $\mathbb{R}^{+}$ 上定义函数 $\Phi$ 如下：
\begin{equation}\tag{24}
\Phi(e^{v}+e^{-v}-2)=\phi\begin{pmatrix}e^{v/2}&\\&e^{-v/2}\end{pmatrix}.
\end{equation}

\par\medskip\noindent\textbf{引理 24.}\quad\emph{若 $U=e^{u}+e^{-u}-2$，则 $g(u)=Q(U)$，其中}
\begin{equation}\tag{25}
Q(U)=\int_{U}^{\infty}\frac{\Phi(V)\,dV}{\sqrt{V-U}} .
\end{equation}

\smallskip
\noindent\emph{证明.}\ 令 $V=e^{v}+e^{-v}-2$。取 $y=e^{u}$、$\eta=e^{v}$，$x$ 如引理 23，
把 $x$ 看作 $V\ge U$ 的函数。则 $dx=\tfrac12 e^{u}(V-U)^{-1/2}dV$。
我们对 $V$ 从 $U$ 到无穷积分，再乘 $2$ 以计入 $x$ 的正负两支。\hfill$\square$

\par\medskip\noindent\textbf{引理 25.}\quad\emph{有}
\begin{equation}\tag{26}
\Phi(U)=-\frac{1}{\pi}\int_{U}^{\infty}\frac{Q'(V)\,dV}{\sqrt{V-U}} .
\end{equation}

\smallskip
\noindent\emph{证明.}\ 我们想在 (25) 的积分号下求导，但左端点依赖 $U$，须当心。
对 (25) 分部积分，得
\[
Q(U)=-2\int_{U}^{\infty}\Phi'(V)\sqrt{V-U}\,dV .
\]
此时被积函数在左端点为零，故可在积分号下求导，再分部积分一次，得
\[
Q'(U)=\int_{U}^{\infty}\frac{\Phi'(V)}{\sqrt{V-U}}\,dV
=-2\int_{U}^{\infty}\Phi''(V)\sqrt{V-U}\,dV .
\]
把它代入 (26) 的右端，再交换积分次序：
\[
\begin{aligned}
-\frac{1}{\pi}\int_{U}^{\infty}\frac{Q'(V)\,dV}{\sqrt{V-U}}
&=\frac{2}{\pi}\int_{U}^{\infty}\int_{V}^{\infty}
\sqrt{\frac{W-V}{V-U}}\,\Phi''(W)\,dW\,dV\\
&=\frac{2}{\pi}\int_{U}^{\infty}\int_{U}^{W}
\sqrt{\frac{W-V}{V-U}}\,dV\,\Phi''(W)\,dW .
\end{aligned}
\tag{27}
\]
为计算内层积分，令 $V=U+(W-U)v$。则
\[
\int_{U}^{W}\sqrt{\frac{W-V}{V-U}}\,dV=(W-U)\int_{0}^{1}\sqrt{\frac{1-v}{v}}\,dv
=\frac{\pi}{2}(W-U),
\]
因为 $\int_{0}^{1}v^{-1/2}(1-v)^{1/2}dv=B(\tfrac12,\tfrac32)=\tfrac{\pi}{2}$。
于是 (27) 等于
\[
\int_{U}^{\infty}(U-W)\Phi''(W)\,dW=-\int_{U}^{\infty}\Phi'(W)\,dW=\Phi(U),
\]
这里用了两次分部积分。\hfill$\square$

\smallskip
\noindent\emph{定理 22 的证明.}\ 取 $U=0$，把引理 25 写成
\begin{equation}\tag{28}
\phi(1)=\Phi(0)=-\frac{1}{\pi}\int_{0}^{\infty}
\frac{g'(u)\,du}{e^{u/2}-e^{-u/2}} .
\end{equation}
$g$ 是偶函数而 $g'$ 是奇函数，故有 Fourier 反演公式
\[
g(u)=\frac{1}{\pi}\int_{0}^{\infty}h(t)\,e^{-itu}\,dt,\qquad
g'(u)=\frac{1}{i\pi}\int_{0}^{\infty}t\,h(t)\,e^{-itu}\,dt,
\]
于是可把 (28) 的积分限改为 $(-\infty,\infty)$ 并除以 $2$，再交换积分次序，得到
\[
\phi(1)=\frac{1}{2\pi}\int_{0}^{\infty}t\,h(t)\,\frac{i}{\pi}
\int_{-\infty}^{\infty}\frac{e^{-iut}\,du}{e^{u/2}-e^{-u/2}}\,dt .
\]
内层积分经过被积函数在 $u=0$ 处的极点，按主值理解。Plancherel 公式随即由
\begin{equation}\tag{29}
\int_{-\infty}^{\infty}\frac{e^{-iut}\,du}{e^{u/2}-e^{-u/2}}=-i\pi\tanh(\pi t)
\end{equation}
推出。此式可如下证明。因 $t>0$，分子 $e^{-iut}$ 在下半平面很小，故可把积分路径向下移。
(29) 的左端等于 $-2\pi i$ 乘以下半平面内被积函数在 $u=-2\pi in$ 处留数之和。
$u=0$ 处的留数只计一半，因为积分路径经过此点。于是得到
\[
-2\pi i\left(\frac12+\sum_{k=1}^{\infty}(-1)^{k}e^{-2\pi kt}\right)=-\pi i\tanh(\pi t),
\]
(i) 得证。

为证 (ii)，定义 $\phi_{2}'(g)=\overline{\phi_{2}(g^{-1})}$。$\phi_{2}$ 是 $K$-双不变的，
而 $K$-双陪集在 $g\to g^{-1}$ 下稳定，故此式等于 $\overline{\phi_{2}(g)}$。
现对 $\phi=\phi_{1}*\phi_{2}'$ 用 (i)，则 $\phi(0)=\langle\phi_{1},\phi_{2}\rangle$。
另一方面 $h_{\phi}(t)=h_{\phi_{1}}(t)\overline{h_{\phi_{2}}(t)}$，由定理 20 即得 (ii)。
\hfill$\square$

\par\medskip\noindent\textbf{命题 26.}\quad\emph{设 $g$ 是 $\mathbb{R}$ 上偶的紧支函数。
则存在 $\phi\in\mathcal{H}^{\circ}$ 使 (20) 成立。}

\smallskip
\noindent\emph{证明.}\ 由 $Q(e^{u}+e^{-u}-2)=g(u)$ 定义 $Q:\mathbb{R}^{+}\to\mathbb{C}$。
再用 (26)，或分部积分并代换 $W=V-U$，定义 $\Phi:\mathbb{R}^{+}\to\mathbb{C}$：
\[
\Phi(U)=\frac{2}{\pi}\int_{0}^{\infty}Q''(W+U)\sqrt{W}\,dW .
\]
可以在积分号下任意次求导，故 $\Phi$ 光滑，在 $U=0$ 处亦然。
现在 $u\to\Phi(e^{u}+e^{-u}-2)$ 是光滑、偶、紧支的函数，
于是存在唯一的光滑紧支 $K$-双不变函数 $\phi$ 满足 (24)。命题随即由引理 24 推出。
\hfill$\square$

\begin{annote}
定理 22 只涉及偶主系列，因为 $K$-不变向量只出现在 $\epsilon=0$ 的 $\pi_{s}^{+}$ 之中；
若对一般的 $\phi\in\mathcal{H}$（不假定 $K$-双不变）写 Plancherel 公式，
还须加上奇主系列的贡献，测度在 $t$ 上的密度将乘以 $2$。
又 (22) 中的 $\int_{0}^{\infty}$ 与 (37) 中的 $\int_{-\infty}^{\infty}$ 相差因子 $2$，
原因是 $h$ 为偶函数。
\par\vspace{0.15em}{\raggedleft 【来源：Knapp, \emph{Representation Theory of Semisimple Groups}, Princeton, 1986, Ch.~XI；Varadarajan, \emph{An Introduction to Harmonic Analysis on Semisimple Lie Groups}, Cambridge, 1999, Ch.~IV（A）】\par}
\end{annote}

\section{Selberg 迹公式}

我们已经证明积分算子 $T_{\phi}$ 是 Hilbert--Schmidt 的，因而紧。还有更多：
它们是\emph{迹类}的。紧算子 $T$ 称为迹类的，若它可分解为两个 Hilbert--Schmidt 算子的复合。
若 $T$ 自伴、特征值为 $\lambda_{i}$，则容易看出迹类性等价于 $\sum|\lambda_{i}|<\infty$。
Lang [30] 中有许多关于迹类算子的有用材料。

取 $f_{i}$ 为 $L^{2}(\Gamma\backslash\mathfrak{H})$ 的一组基，其中每个 $f_{i}$ 既是 $T_{\phi}$
的特征函数，又是 $\Delta$ 的特征函数。设 $\phi\in\mathcal{H}^{\circ}$ 满足 (6)，以 $\mu_{i}$
记 $T_{\phi}$ 在这组基下的特征值。作 Fourier 展开，得
\begin{equation}\tag{30}
K_{\phi}(z,w)=\sum_{i}\mu_{i}f_{i}(z)\overline{f_{i}(w)} .
\end{equation}

\par\medskip\noindent\textbf{定理 27.}\quad\emph{若 $\phi\in\mathcal{H}^{\circ}$，
则 $T_{\phi}$ 是迹类的。}

\smallskip
\noindent\emph{证明.}\ 迹类算子的线性组合仍是迹类的。故只需对
$\tfrac12(\phi(g)+\overline{\phi(g^{-1})})$ 与 $\tfrac1{2i}(\phi(g)-\overline{\phi(g^{-1})})$
证明结论。于是可设 $\phi$ 满足 (6)，从而 $T_{\phi}$ 自伴。设 $\mu_{i}$ 是它的（非零）特征值，
$\lambda_{i}$ 是 $\Delta$ 的相应特征值。于是 $\sum\lambda_{i}^{-2}<\infty$。

对 $K_{\phi}(z,w)$ 的第一个变量作用 $\Delta$，得到新的核 $\Delta_{z}K_{\phi}$。我们要证
\begin{equation}\tag{31}
(\Delta_{z}K_{\phi})(z,w)=\sum_{i}\mu_{i}\lambda_{i}f_{i}(x)\overline{f_{i}(y)} .
\end{equation}
形式上，由 (30) 逐项求导即得。但 (30) 的收敛是 $L^{2}$ 意义下的，而 $\Delta$ 是无界算子，
故须验证。注意 $\Delta_{x}K_{\phi}(x,y)$ 连续，故有展开
\[
(\Delta_{z}K_{\phi})(z,w)=\sum_{i}\nu_{i}f_{i}(x)\overline{f_{i}(y)} .
\]
设 $s>0$，考虑
\[
\int_{\Gamma\backslash\mathfrak{H}}G_{s}(z,\zeta)
\bigl[(\Delta_{\zeta}K_{\phi})(\zeta,w)+sK_{\phi}(\zeta,w)\bigr]\,
\frac{d\xi\wedge d\eta}{\eta^{2}} .
\]
一方面，由 (9)，这正是 (30)。另一方面，方括号中的项等于
$\sum_{i}(\nu_{i}+s\mu_{i})f_{i}(\zeta)\overline{f_{i}(w)}$；又预解式紧，故可逐项作用。
用 $R(s,\Delta)f_{i}=(s+\lambda_{i})^{-1}f_{i}$ 得 $(\nu_{i}+s\mu_{i})(\lambda_{i}+s)^{-1}=\mu_{i}$，
即 $\nu_{i}=\lambda_{i}\mu_{i}$，(31) 得证。

由于 $\Delta_{z}K(z,w)$ 连续，它是 Hilbert--Schmidt 的，于是得到界
\begin{equation}\tag{32}
\sum|\mu_{i}\lambda_{i}|^{2}<\infty .
\end{equation}
现在由 $\sum|\lambda_{i}|^{-2}<\infty$ 与 (32)，用 Cauchy--Schwarz 即得
$\sum|\mu_{i}|<\infty$。\hfill$\square$

自伴迹类算子 $T$ 的迹 $\operatorname{tr}T$ 按定义是其特征值之和。

\par\medskip\noindent\textbf{定理 28.}\quad\emph{若 $\phi\in\mathcal{H}^{\circ}$ 满足 (6)，
$\mu_{i}$ 是 $T_{\phi}$ 的特征值，则}
\begin{equation}\tag{33}
\operatorname{tr}T_{\phi}=\int_{\Gamma\backslash\mathfrak{H}}K_{\phi}(z,z)\,
\frac{dx\wedge dy}{y^{2}} .
\end{equation}

\smallskip
\noindent\emph{证明.}\ 对 (30) 积分并用标准正交性即得。\hfill$\square$

Selberg 迹公式是迹的一个更显式的公式。以 $\{\gamma\}$ 记 $\Gamma$ 的共轭类的一组代表元，
以 $Z_{\Gamma}(\gamma)$ 记 $\gamma$ 在 $\Gamma$ 中的中心化子。

\par\medskip\noindent\textbf{定理 29.}\quad\emph{有}
\begin{equation}\tag{34}
\operatorname{tr}T_{\phi}=\sum_{\{\gamma\}}\int_{Z_{\Gamma}(\gamma)\backslash G}
\phi(g^{-1}\gamma g)\,dg .
\end{equation}

\smallskip
\noindent\emph{证明.}\ 把 (33) 的右端改写为
\[
\sum_{\gamma\in\Gamma}\int_{G}\phi(g^{-1}\gamma g)\,dg
=\sum_{\{\gamma\}}\sum_{\delta\in Z_{\Gamma}(\gamma)\backslash\Gamma}
\int_{G}\phi(g^{-1}\delta^{-1}\gamma\delta g)\,dg .
\]
把积分与求和合并即得 (34)。\hfill$\square$

这是迹公式的原始形式。要使它更显式，须更具体地研究右端的\emph{轨道积分}。
元素 $1\neq\gamma\in\Gamma$ 称为\emph{双曲的}，若其特征值为实数；
称为\emph{椭圆的}，若特征值为模 $1$ 的复数。若每个 $1\neq\gamma\in\Gamma$ 都是双曲的，
则称 $\Gamma$ 是双曲的。例如取 $X$ 为亏格 $\ge2$ 的紧 Riemann 曲面，
其万有覆叠是 $\mathfrak{H}$，而 $\Gamma=\pi_{1}(X)$ 以 $X$ 为商。这些例子恰是双曲群。

以下设 $\Gamma$ 是双曲群。若 $1\neq\gamma\in\Gamma$，定义 $N=N(\gamma)$ 为使 $\gamma$
共轭于
\begin{equation}\tag{35}
\begin{pmatrix}N^{1/2}&\\&N^{-1/2}\end{pmatrix}
\end{equation}
的数。再令 $N_{0}=N_{0}(\gamma)$ 使
\[
\begin{pmatrix}N_{0}^{1/2}&\\&N_{0}^{-1/2}\end{pmatrix}
\]
共轭于 $Z_{\Gamma}(\gamma)$ 的一个生成元。显然可设 $N$ 与 $N_{0}$ 都大于 $1$。
注意 $Z_{G}(\gamma)$ 共轭于对角子群，它在 $X$ 中的像是闭测地线。
于是诸 $N(\gamma)$ 是闭测地线的长度，诸 $N_{0}(\gamma)$ 是素测地线的长度。

\par\medskip\noindent\textbf{定理 30.}\quad\emph{取 $g$ 为 (20) 中的函数，则}
\begin{equation}\tag{36}
\int_{Z_{\Gamma}(\gamma)\backslash G}\phi(g^{-1}\gamma g)\,dg
=\frac{\log N_{0}}{N^{1/2}-N^{-1/2}}\,g(\log N).
\end{equation}

\smallskip
\noindent\emph{证明.}\ 可设 $\gamma$ 等于 (35)。用 Iwasawa 坐标 (23)，该积分为
\[
\begin{aligned}
&\int_{0}^{\log N_{0}}du\int_{-\infty}^{\infty}
\phi\left(\begin{pmatrix}1&-x\\&1\end{pmatrix}
\begin{pmatrix}N^{1/2}&\\&N^{-1/2}\end{pmatrix}
\begin{pmatrix}1&x\\&1\end{pmatrix}\right)dx\\
&\quad=|\log(N_{0})|\int_{-\infty}^{\infty}
\phi\left(\begin{pmatrix}N^{1/2}&\\&N^{-1/2}\end{pmatrix}
\begin{pmatrix}1&(1-N^{-1})x\\&1\end{pmatrix}\right)dx,
\end{aligned}
\]
再作变量替换即得 (36)。\hfill$\square$

\par\medskip\noindent\textbf{定理 31.}\quad\emph{设 $g$ 是光滑、偶、紧支的函数，
$h$ 是它由 (21) 定义的 Fourier 变换。若 $\tfrac14+t_{i}^{2}$ 是 $\Delta$ 在 $\mathfrak{H}$ 上的
特征值，而 $\log(N)$ 跑遍 $\Gamma$ 的闭测地线长度，其中对每个 $N$ 以 $N_{0}$ 记相应素测地线的长度，
则}
\begin{equation}\tag{37}
\sum h(t_{i})=\frac{\operatorname{vol}(\Gamma\backslash\mathfrak{H})}{4\pi}
\int_{-\infty}^{\infty}t\,h(t)\tanh(\pi t)\,dt
+\sum_{N}\frac{\log N_{0}}{N^{1/2}-N^{-1/2}}\,g(\log N).
\end{equation}

\smallskip
这就是 \emph{Selberg 迹公式}。

\smallskip
\noindent\emph{证明.}\ 按命题 26 选取 $\phi$。由定理 19 与定理 21，诸 $h(t_{i})$ 是
$T_{\phi}$ 在 Laplacian 的特征函数上的特征值，故 (37) 的左端是 $\operatorname{tr}T_{\phi}$。
由定理 22 与定理 30，(37) 的右端是诸轨道积分之和。于是恒等式由定理 29 推出。
\hfill$\square$

\begin{annote}
(37) 右端两项各有几何含义。第一项是恒等项，只依赖 $\Gamma\backslash\mathfrak{H}$ 的体积；
取 $h$ 为近似 $\delta$ 的函数即得 Weyl 律（定理 33）。第二项是双曲项，
按素测地线的长度 $\log N_{0}$ 求和，系数 $\log N_{0}/(N^{1/2}-N^{-1/2})$ 是测地流在该轨道上的周期。
一般的 $\Gamma$ 还须补上椭圆共轭类的贡献，即所谓椭圆项；原书为简化只讨论双曲群，
故 (37) 中没有这一项。
\par\vspace{0.15em}{\raggedleft 【来源：Hejhal, Vol.~1, Springer LNM 548, 1976, Introduction；Iwaniec, 1995, Ch.~10（A）】\par}
\end{annote}

\section{Selberg zeta 函数}

为得到有用的应用，必须扩大迹公式中函数 $g$ 与 $h$ 的类。

\par\medskip\noindent\textbf{定理 32.}\quad\emph{定理 31 的迹公式仍然成立，只要 $h$ 是偶函数，
在带形 $|\operatorname{im}z|\le\tfrac12+\delta$ 内解析，且在此带形内满足
$h(r)=O((1+|r|)^{-2-\delta})$。这一假设蕴涵 Fourier 变换满足
$g(u)=O(e^{-(\frac12+\delta)|u|})$。}

\smallskip
用逼近论证由定理 31 推出这个更强陈述的证明，见 Hejhal [20, Chapter 1, Section 7]。
\hfill$\square$

\par\medskip\noindent\textbf{定理 33（Weyl 律）.}\quad\emph{满足 $t_{j}\le x$ 的 $j$ 的个数
渐近于 $\frac{1}{4\pi}\operatorname{vol}(\Gamma\backslash\mathfrak{H})x^{2}$。}

\smallskip
这里必须提到，这是迹公式的第一个重要应用。注意 $\lambda_{j}=\tfrac14+t_{j}^{2}$，
故满足 $\lambda_{j}\le x$ 的 $j$ 的个数渐近于
$\frac{1}{4\pi}\operatorname{vol}(\Gamma\backslash\mathfrak{H})x$。

我们不证明 Weyl 律。简述如下：取 $h(t)=e^{-t^{2}T}$，则定理 31 中双曲项的量级较小，
(37) 的第一项占优。于是得到 $\sum e^{-\lambda_{j}T}\cong\frac{1}{4\pi}
\operatorname{vol}(\Gamma\backslash\mathfrak{H})$，再由 Tauber 定理即得 Weyl 律。
见 Hejhal [20, Chapter 2, Section 2]。\hfill$\square$

迹公式可以看作 $\Gamma\backslash\mathfrak{H}$ 的长度谱（即闭测地线长度 $N$ 之集）
与那些使 $\tfrac14+t_{i}^{2}$ 为 Laplacian 特征值的数 $t_{i}$ 之间的对偶。
下面讨论的另一种类似的对偶，存在于素数集与 Riemann zeta 函数的零点之间。

\par\medskip\noindent\textbf{定理 34.}\quad\emph{设 $g$ 紧支、光滑且偶，
$h(t)=\int_{-\infty}^{\infty}e^{itu}g(u)\,du$。以 $\tfrac12+it_{i}$ 记 $\zeta$
在临界带内的零点。则}
\begin{equation}\tag{38}
\begin{aligned}
h\!\left(-\frac{i}{2}\right)-\sum h(t_{i})+h\!\left(\frac{i}{2}\right)
&=-\frac{1}{2\pi}\int_{-\infty}^{\infty}h(t)\,
\frac{\Gamma'}{\Gamma}\!\left(\frac{1+2it}{4}\right)dt\\
&\quad+g(0)\log(\pi)+2\sum\frac{\log(p)}{\sqrt{p^{n}}}\,g(\log(p^{n})).
\end{aligned}
\end{equation}

\smallskip
既然假定 $h$ 是偶函数，就有 $h(-i/2)=h(i/2)$。然而把公式写成这样是有道理的，
见下面的注 4。特殊情形（“显式公式”）由 Riemann、von Mangoldt、Hadamard、
de la Vallée Poussin 与 Ingham 得到，Ingham [24] 中有讨论。Weil [50] 把这层对偶形式化了。

\begin{annote}
(38) 即 Riemann--von Mangoldt 显式公式。左端 $h(-i/2)$ 与 $h(i/2)$ 来自 $\xi(s)$ 在
$s=0,1$ 处的平凡零点，中间的和遍历非平凡零点 $\tfrac12+it_{i}$；右端第一项是 archimedean 因子
$\Gamma_{\mathbb{R}}$ 的对数导数，末项由 $-\zeta'/\zeta=\sum\Lambda(n)n^{-s}$ 的素数幂给出。
原书把 $h(-i/2)$ 与 $h(i/2)$ 分开写并非多余：它们分别对应 $H^{0}$ 与 $H^{2}$ 的贡献，
注 4 的有限域类比正是三部分迹之和。
\par\vspace{0.15em}{\raggedleft 【来源：Ingham, \emph{The Distribution of Prime Numbers}, Cambridge, 1932, Ch.~IV；Weil, \emph{Sur les ``formules explicites'' de la th\'eorie des nombres premiers}, 1952（A）】\par}
\end{annote}


\smallskip
\noindent\emph{证明.}\ 为此令 $\xi(s)=\Gamma_{\mathbb{R}}(s)\zeta(s)$，其中
$\Gamma_{\mathbb{R}}(s)=\pi^{-s/2}\Gamma\left(\frac{s}{2}\right)$。
函数 $h(t)$ 整，我们写 $h(t)=H\left(\frac12+it\right)$，$H(s)=h(-i(s-\frac12))$。考虑
\[
\frac{1}{2\pi i}\int_{1+\delta-i\infty}^{1+\delta+i\infty}\frac{\xi'(s)}{\xi(s)}H(s)\,ds,
\qquad\delta>0 .
\]
把积分路径移到 $\operatorname{re}(s)=-\delta$，并用函数方程 $\xi(s)=\xi(1-s)$，
得到上述积分的相反数再加上留数之和，于是
\[
\frac{1}{\pi i}\int_{1+\delta-i\infty}^{1+\delta+i\infty}\frac{\xi'(s)}{\xi(s)}H(s)\,ds
=\sum h(t_{i})-h\!\left(\frac{i}{2}\right)-h\!\left(-\frac{i}{2}\right).
\]
又有 $\xi'/\xi=\Gamma_{\mathbb{R}}'/\Gamma_{\mathbb{R}}+\zeta'/\zeta$，
而关于 $\Gamma_{\mathbb{R}}$ 的积分可以向左移到 $\operatorname{re}(s)=\tfrac12$。于是
\[
\begin{aligned}
&\frac{1}{\pi i}\int_{\frac12-i\infty}^{\frac12+i\infty}
\frac{\Gamma_{\mathbb{R}}'(s)}{\Gamma_{\mathbb{R}}(s)}\,h\!\left(-i\left(s-\frac12\right)\right)ds\\
&\qquad=-g(0)\log(\pi)+\frac{1}{2\pi}\int_{-\infty}^{\infty}h(t)\,
\frac{\Gamma'}{\Gamma}\!\left(\frac{1+2it}{4}\right)dt .
\end{aligned}
\]
另一方面，$-\zeta'(s)/\zeta(s)=\sum\Lambda(n)n^{-s}$，其中对素数 $p$ 有
$\Lambda(p^{k})=\log(p)$，而 $n$ 不是素数幂时 $\Lambda(n)=0$。又有
\[
\frac{1}{2\pi i}\int_{2-i\infty}^{2+i\infty}H(s)n^{-s}\,ds=\frac{g(\log(n))}{\sqrt{n}},
\]
把各部分拼起来即得 (38)。\hfill$\square$

\par\medskip\noindent\textbf{注 1.}\ 移动积分线时须知道：对任意大的 $T$，都能找到从 $2+iT$
到 $-1+iT$ 的路径，使 $\zeta'/\zeta$ 在其上不太大。办法是让路径大致走在一对相距不太近的零点
中间。见 Ingham [24, Theorem 2, p.\ 71]，那里证明总能找到 $T$ 附近的 $t$，
使 $\zeta'/\zeta$ 在从 $2+it$ 到 $-1+it$ 的直线上为 $O(\log^{2}(T))$。

\par\medskip\noindent\textbf{注 2.}\ 实践中 $g$ 紧支这个条件太强。只需 $g$ 与 $h$
如定理 32 所述即可。

\par\medskip\noindent\textbf{注 3.}\ 若 $H$ 在左半平面不太大，则有
\[
\frac{1}{2\pi}\int_{-\infty}^{\infty}h(t)\,
\frac{\Gamma'\!\left(\frac{1+2it}{4}\right)}{\Gamma\!\left(\frac{1+2it}{4}\right)}\,dt
=\frac{1}{2\pi i}\int_{\frac12-i\infty}^{\frac12+i\infty}H(s)\,
\frac{\Gamma'\!\left(\frac{s}{2}\right)}{\Gamma\!\left(\frac{s}{2}\right)}\,ds
=2\sum_{n=0}^{\infty}H(-2n),
\]
其中求和跑遍 $\Gamma(s/2)$ 的极点处的留数。

\par\medskip\noindent\textbf{注 4.}\ 设 $X$ 是有限域 $k=\mathbb{F}_{q}$ 上亏格为 $g$ 的
非奇异完全曲线。若 $P$ 是次数为 $d(P)$ 的素除子，记 $N(P)=q^{d(P)}$。$X$ 的 zeta 函数是
\[
\prod_{P}(1-N(P)^{-s})^{-1}
=\frac{\prod_{j=1}^{2g}(1-\alpha_{j}q^{-s})^{-1}}{(1-q^{-s})(1-q^{1-s})},
\]
其中 $\alpha_{j}$ 是 Frobenius 映射在 $H^{1}(X,\mathbb{Q}_{l})$ 上的特征值，
$l\neq p$ 为任一素数。Riemann 假设即 $|\alpha_{j}|=\sqrt{q}$。
写 $\alpha_{j}=q^{\frac12+i\theta_{j}}$。设 $g(n):\mathbb{Z}\to\mathbb{C}$ 是偶的序列，
且只有有限多个 $n$ 使其非零，并令
\[
h(t)=\sum_{n}g(n)\,e^{int\log(q)} .
\]
于是 $h$ 整、偶，且以 $2\pi/\log(q)$ 为周期。此时
\[
h\!\left(-\frac{i}{2}\right)-\sum h(\theta_{j})+h\!\left(\frac{i}{2}\right)
=-(2g-2)g(0)+\sum_{P}\sum_{m=1}^{\infty}
\frac{d(P)\,g(m\,d(P))}{\sqrt{q^{md(P)}}} .
\]
定理 34 的这个函数域类似物可通过考虑
\[
\frac{1}{2\pi i}\int_{2}^{2+2\pi i/\log(q)}
\left(\frac{Z'(s)}{Z(s)}+(g-1)\log(q)\right)H(s)\,ds
\]
来证明，其中 $h(t)=H\left(\frac12+it\right)$。这可以看作 Lefschetz--Grothendieck
不动点公式应用于 $X$ 上一个对应所得的结果，而 $h(-i/2)$、$\sum h(\theta_{j})$、
$h(i/2)$ 三项分别是该对应作用在 $H^{0}(X)$、$H^{1}(X)$、$H^{2}(X)$ 上的迹。
见 Patterson [38, Chapter 5] 与 Connes [10] 的第一部分。

Selberg 迹公式与显式公式之间的类似，是许多推测的来源。这种类似是否会引起误解，
尚待分晓。显式公式与 Selberg 迹公式的一处重大差别在于：迹公式中 $\sum h(t_{i})$
以正号出现，而在 (38) 中以负号出现。正如 Connes [10] 所指出的，这一差别可能非常重要。

这种类似的一项切实成果，是 Selberg 发现了一个 zeta 函数，它与迹公式的关系，
如同 Riemann zeta 函数与显式公式的关系。沿用上一节的记号，Selberg 考虑
\[
Z(s)=\prod_{\{N_{0}\}}\prod_{k=0}^{\infty}(1-N_{0}^{-s-k}),
\]
其中 $\log(N_{0})$ 跑遍素测地线的长度。

\par\medskip\noindent\textbf{定理 35.}\quad\emph{$Z(s)$ 解析延拓到一切 $s$，
其零点在负整数处，以及在 $\tfrac12\pm it_{k}$ 处，
其中 $\tfrac14+t_{k}^{2}$ 是 Laplacian 在 $\Gamma\backslash\mathfrak{H}$ 上的特征值。}

\smallskip
\noindent\emph{证明与讨论.}\ 为说明引入 Selberg zeta 函数的动机，我们本想在迹公式中取
\[
g(u)=e^{-|u|(s-1/2)} .
\]
可惜不能使用这个函数。但若能使用，迹公式的几何侧就会是 $Z(s)$ 的对数导数。

\par\medskip\noindent\textbf{引理 36.}\quad
\[
\sum_{\{N\}}\frac{\log(N_{0})\,N^{-(s-1/2)}}{N^{1/2}-N^{-1/2}}=\frac{Z'(s)}{Z(s)} .
\]

\smallskip
\noindent\emph{证明.}\ 写 $N=N_{0}^{m}$，则
\[
\sum_{\{N_{0}\}}\sum_{m=1}^{\infty}
\frac{\log(N_{0})\,N_{0}^{-m(s-1/2)}}{N_{0}^{m/2}-N_{0}^{-m/2}}
=\sum_{\{N_{0}\}}\sum_{m=1}^{\infty}
\frac{\log(N_{0})\,N_{0}^{-ms}}{1-N_{0}^{-m}} .
\]
代入 $(1-N_{0}^{-m})^{-1}=\sum_{k=0}^{\infty}N_{0}^{-ks}$ 并交换求和次序，此式等于
\[
\sum_{\{N_{0}\}}\frac{\log(N_{0})\,N_{0}^{-(s+k)}}{1-N_{0}^{-(s+k)}}=\frac{Z'(s)}{Z(s)} .
\]
证毕。\hfill$\square$

虽然不能取 $g(u)=e^{-|u|(s-1/2)}$，但可以取
\[
g(u)=\frac{e^{-|u|(s-1/2)}}{2s-1}-\frac{e^{-|u|(\sigma-1/2)}}{2\sigma-1},
\]
其中 $s$ 与 $\sigma$ 不同。这一相减的技巧消去了 $g'$ 在 $u=0$ 处的不连续性。
我们固定 $\sigma$ 而变动 $s$。起初假定两者的实部都很大。Fourier 变换
$h(t)=\int_{-\infty}^{\infty}g(u)e^{itu}du$ 为
\[
h(t)=\frac{1}{\left(s-\frac12\right)^{2}+t^{2}}
-\frac{1}{\left(\sigma-\frac12\right)^{2}+t^{2}} .
\]
其次我们证明：设 $s$ 与 $\sigma$ 的实部都 $\ge\tfrac12$，则
\begin{equation}\tag{39}
\int_{-\infty}^{\infty}h(t)\,t\tanh(\pi t)\,dt
=2\sum_{k=0}^{\infty}\left[\frac{1}{s+k}-\frac{1}{\sigma+k}\right].
\end{equation}
为此用部分分式分解
\[
\frac{t}{\left(s-\frac12\right)^{2}+t^{2}}
=\frac{1}{2i}\left[\frac{1}{\left(s-\frac12\right)-it}
-\frac{1}{\left(s-\frac12\right)+it}\right].
\]
(39) 的左端等于
\[
\begin{aligned}
&\frac{1}{2i}\int_{-\infty}^{\infty}\left[
\frac{1}{\left(s-\frac12\right)-it}-\frac{1}{\left(\sigma-\frac12\right)-it}\right]
\tanh(\pi t)\,dt\\
&\quad-\frac{1}{2i}\int_{-\infty}^{\infty}\left[
\frac{1}{\left(s-\frac12\right)+it}-\frac{1}{\left(\sigma-\frac12\right)+it}\right]
\tanh(\pi t)\,dt .
\end{aligned}
\]
由于 $\tanh$ 是奇函数，两部分贡献相等，而第一个积分可把路径向下移到下半平面来计算，
那里只有 $\tanh(\pi t)$ 的极点 $-i\left(k+\frac12\right)$，$k=0,1,2,\dots$。
$\tanh(\pi t)$ 在这些点的留数是 $\pi^{-1}$，于是得到 (39)。

现在使用 Selberg 迹公式。注意基本区域的体积是 $4\pi(g-1)$。
于是对这些函数 $g$ 与 $h$，我们得到
\[
\begin{aligned}
&\frac{1}{2s-1}\frac{Z'(s)}{Z(s)}-\frac{1}{2\sigma-1}\frac{Z'(\sigma)}{Z(\sigma)}\\
&\quad=-(2g-2)\sum_{k=0}^{\infty}\left[\frac{1}{s+k}-\frac{1}{\sigma+k}\right]
+\sum_{i}\left[\frac{1}{\left(s-\frac12\right)^{2}+t_{i}^{2}}
-\frac{1}{\left(\sigma-\frac12\right)^{2}+t_{i}^{2}}\right].
\end{aligned}
\]
固定 $\sigma$ 而变动 $s$，右端亚纯延拓到一切 $s$。乘以 $2s-1$ 可见 $Z'/Z$ 的极点是单极点，
留数是整数。于是
\[
Z(s)=\exp\int_{\infty}^{s}\frac{Z'(s)}{Z(s)}\,ds
\]
解析延拓到一切 $s$。\hfill$\square$

\begin{annote}
由定理 35，$Z(s)$ 的零点有两类：$s=-k$（$k=0,1,2,\dots$，来自无穷乘积中各因子的平凡零点）
与 $s=\tfrac12\pm it_{k}$，其中 $\lambda_{k}=\tfrac14+t_{k}^{2}$ 遍历 $\Delta$ 的特征值。
故 Riemann 假设对 $Z$ 而言无条件成立；真正的困难在例外特征值：若 $t_{k}$ 为纯虚数，
对应零点就落在实轴 $(\tfrac12,1)$ 内，Selberg 猜想正是要排除这种情形。
\par\vspace{0.15em}{\raggedleft 【来源：Selberg, J. Indian Math. Soc. 20 (1956)；Hejhal, Vol.~2, Ch.~2（A）】\par}
\end{annote}

\section{带一个尖点的群}

若群 $\Gamma$ 的体积有限但带尖点，则谱理论因连续谱而变得复杂，连续谱来自 Eisenstein 级数。
不过尖点谱仍是离散的。为简单起见，设 $\infty$ 是 $\Gamma$ 的唯一尖点，
且 $\infty$ 在 $\Gamma$ 中的稳定化子是
\[
\Gamma_{\infty}=\left\{\begin{pmatrix}1&n\\&1\end{pmatrix}\Bigm|n\in\mathbb{Z}\right\}.
\]
我们还将记 $G_{\infty}=\left\{\begin{pmatrix}1&x\\&1\end{pmatrix}\Bigm|x\in\mathbb{R}\right\}$。
以 $L_{0}^{2}(\Gamma\backslash G)$ 记由“尖点形式” $f\in L^{2}(\Gamma\backslash G)$ 构成的子空间，
它们满足
\begin{equation}\tag{40}
\int_{0}^{1}f\left(\begin{pmatrix}1&x\\&1\end{pmatrix}g\right)dx=0 .
\end{equation}
再以 $L_{0}^{2}(\Gamma\backslash\mathfrak{H})$ 记 $L_{0}^{2}(\Gamma\backslash G)$ 中右 $K$-不变的
元素，它们可以看作 $\mathfrak{H}$ 上的函数。

与定理 7 相反，算子 $T_{\phi}$ 不再紧。但我们仍有：

\par\medskip\noindent\textbf{定理 37（Gelfand--Graev--Piatetski-Shapiro）.}\quad
\emph{若 $\phi\in\mathcal{H}$，则 $T_{\phi}$ 在 $L_{0}^{2}(\Gamma\backslash G)$ 上的限制是紧算子。}

\smallskip
\noindent\emph{证明.}\ Godement 的证明见 Bump [6, Section 3.2]。\hfill$\square$

\par\medskip\noindent\textbf{定理 38.}\quad
(i) \emph{$L_{0}^{2}(\Gamma\backslash\mathfrak{H})$ 有一组由 $\Delta$ 的特征函数构成的基。}

(ii) \emph{$L_{0}^{2}(\Gamma\backslash G)$ 分解为不可约不变子空间的直和。
这些子空间中任一空间的 $K$-有限元素都速降。}

\smallskip
\noindent\emph{证明梗概.}\ 定理 9 与定理 11 的证明稍作修改即可。
对 (ii) 的第二个断言，只需证明：若 $V\subset L_{0}^{2}(\Gamma\backslash G)$ 不可约、
$f\in V(k)$，则 $f$ 速降。我们可以找到 $\phi\in\mathcal{H}$，满足
$\phi(\kappa_{\theta}g\kappa_{\sigma})=e^{ik(\theta+\sigma)}\phi(g)$，
并使 $y\mapsto\phi\begin{pmatrix}y^{1/2}&\\&y^{-1/2}\end{pmatrix}$ 是总量为 $1$、
集中在 $y=1$ 附近的正函数。于是 $T_{\phi}f$ 接近 $f$，因而非零；而它在 $V(k)$ 中，
由命题 5 后者一维，故 $T_{\phi}f$ 与 $f$ 成比例。由核 $K_{0}(g,h)$ 的速降性
（定义见 Bump [6, Proposition 3.2.3]）知 $T_{\phi}f$ 速降，从而 $f$ 速降。
\hfill$\square$

与紧商的情形一样，Laplacian 在每个不可约一维子空间上作用为纯量。
尖点谱的性状与紧商情形的整个谱十分相似。另一方面，$L_{0}^{2}(\Gamma\backslash G)$
的正交补含连续谱。与谱定理相关的 Laplacian 特征函数本身并不平方可积。

为理解这是怎么回事，并对连续谱的性质获得一些直观，看下面的例子。
群 $\mathbb{R}$ 通过平移作用在自己上，Laplacian $-d^{2}/dx^{2}$ 是不变微分算子。
它有特征函数 $f_{a}(x)=e^{2\pi iax}$，特征值为 $a^{2}$。任何 $L^{2}$ 函数都有 Fourier 展开
\[
\phi(x)=\int_{-\infty}^{\infty}\hat\phi(a)f_{a}(x)\,da,
\]
但 $f_{a}$ 本身不属于 $L^{2}$。若 $T\subset\mathbb{R}$ 可测，则支在 $T$ 上的 $L^{2}$ 函数的
Fourier 变换构成不变子空间。这里没有极小不变子空间，
故 $L^{2}(\mathbb{R})$ 不能分解为不可约表示的直和。

为了陈述关于 Eisenstein 级数的基本定理（这些级数既属于奇主系列又属于偶主系列），
我们把设置稍加推广。设离散群 $\Gamma$ 含 $-I$，它在 $\mathfrak{H}$ 上作用平凡。
固定 $\Gamma$ 的一个酉特征 $\chi$，使 $\chi(-I)\equiv(-1)^{\epsilon}\ \mathrm{mod}\ 2$，
且 $\chi|_{\Gamma_{\infty}}$ 平凡，其中 $\epsilon=0$ 或 $1$。
令 $k\equiv\epsilon$ 模 $2$，并令 $\pm=(-1)^{\epsilon}$。回忆 $f_{s,k}$ 由 (15) 定义，
\begin{equation}\tag{41}
E(g,f_{s,k},\chi)=\sum_{\Gamma_{\infty}\backslash\Gamma}
\overline{\chi(\gamma)}\,f_{s,k}(\gamma g).
\end{equation}

\par\medskip\noindent\textbf{命题 39.}\quad\emph{若 $\operatorname{re}(s)>1$，级数 (41)
绝对收敛。}

\smallskip
\noindent\emph{证明梗概.}\ 设 $\sigma$ 是 $s$ 的实部。要证的是：若 $\sigma>1$，则
\begin{equation}\tag{42}
\sum_{\Gamma_{\infty}\backslash\Gamma}f_{\sigma,0}(\gamma g)<\infty .
\end{equation}
在上半平面上定义测度 $\mu_{\sigma}=y^{\sigma-2}dx\wedge dy$。
设 $B$ 是 $z\in\mathfrak{H}$ 附近的小邻域。映射 $z\to\gamma(z)$ 的 Jacobi 行列式是
$|cz+d|^{-4}$，其中 $\gamma=\binom{a\ b}{c\ d}\in\mathrm{SL}(2,\mathbb{R})$。
于是 $\gamma(B)$ 的 $\mu_{\sigma}$-体积约为 $|cz+d|^{-2\sigma}\operatorname{vol}(B)$。

我们可以选取代表元 $\gamma\in\Gamma_{\infty}\backslash\Gamma$，
使诸像 $\gamma(B)$ 都落在矩形 $0\le x\le1$、$0<y\le C$ 内，$C$ 为某常数。
（其实这并不完全正确。若某个 $\gamma(B)$ 恰好落在此区域的左边界或右边界上，
就沿该边界把它切成两块，并把其中一块平移 $\pm1$ 移回区域内。）
此矩形体积有限，故 $\sum_{\Gamma_{\infty}\backslash\Gamma}|cz+d|^{-2\sigma}<\infty$，
从而 (42) 成立。\hfill$\square$

Eisenstein 级数的解析延拓与它的\emph{常数项}的解析延拓密切相关：
\begin{equation}\tag{43}
E_{0}(g,f_{s,k},\chi)=\int_{0}^{1}
E\left(\begin{pmatrix}1&t\\&1\end{pmatrix}g,\,f_{s,k},\,\chi\right)dt .
\end{equation}
回忆由 (16) 定义的缠绕算子 $M(s):P_{s}^{\pm}\to P_{1-s}^{\pm}$。

\par\medskip\noindent\textbf{命题 40.}\quad\emph{设 $\operatorname{re}(s)>1$。
存在解析函数 $c(s)$，它与 $k$ 无关，在竖带（$1$ 的右侧）上有界，使}
\begin{equation}\tag{44}
E_{0}(g,f_{s,k},\chi)=f_{s,k}(g)+c(s)\,M(s)f_{s,k}.
\end{equation}

\smallskip
\noindent\emph{证明.}\ 把 $E(g,f_{s,k},\chi)$ 的定义代入 (43)。陪集 $\Gamma_{\infty}$
的贡献是 $f_{s,k}$。其余项贡献
\[
\begin{aligned}
&\int_{0}^{1}\sum_{\substack{\gamma\in\Gamma_{\infty}\backslash\Gamma/\Gamma_{\infty}\\
\gamma\notin\Gamma_{\infty}}}\ \sum_{\delta\in\Gamma_{\infty}}
\overline{\chi(\gamma\delta)}\,f_{s,k}\left(\gamma\delta
\begin{pmatrix}1&x\\&1\end{pmatrix}g\right)dx\\
&\qquad=\sum_{\substack{\gamma\in\Gamma_{\infty}\backslash\Gamma/\Gamma_{\infty}\\
\gamma\notin\Gamma_{\infty}}}\overline{\chi(\gamma)}
\int_{-\infty}^{\infty}f_{s,k}\left(\gamma
\begin{pmatrix}1&x\\&1\end{pmatrix}g\right)dx .
\end{aligned}
\]
若 $\gamma=\binom{a\ b}{c\ d}\notin\Gamma_{\infty}$，则 $c\neq0$，且
\[
\gamma=\begin{pmatrix}c^{-1}&a\\&c\end{pmatrix}
\begin{pmatrix}&-1\\1&\end{pmatrix}
\begin{pmatrix}1&d/c\\&1\end{pmatrix},
\]
故变量替换 $x\to x-d/c$ 表明
\[
\int_{-\infty}^{\infty}f_{s,k}\left(\gamma
\begin{pmatrix}1&x\\&1\end{pmatrix}g\right)dx=|c|^{-2s}M(s)f_{s,k}.
\]
于是 (44) 成立，其中
\begin{equation}\tag{45}
c(s)=\sum_{\substack{\gamma\in\Gamma_{\infty}\backslash\Gamma/\Gamma_{\infty}\\
\gamma\notin\Gamma_{\infty}}}\chi(\gamma)\,|c|^{-2s}.
\end{equation}
由于我们处在 Eisenstein 级数绝对收敛的区域内，这个 Dirichlet 级数当 $\operatorname{re}(s)>1$
时收敛；(45) 表明它是竖带有界的解析函数。\hfill$\square$

可以把此式改写成
\begin{equation}\tag{46}
E_{0}(g,f_{s,k},\chi)=f_{s,k}(g)+A_{k}(s)f_{1-s,k},
\end{equation}
其中 $A_{k}$ 等于 $c$ 乘以 (18) 右端的常数。由 Stirling 公式与 $c(\sigma+it)$ 的有界性可得：
当 $\sigma>1$ 时
\[
A(\sigma+it)=O(|t|^{-1/2}).
\]

\begin{annote}
对 $\Gamma=\mathrm{SL}(2,\mathbb{Z})$、$\chi=1$ 有 $A(s)=\xi(2s-1)/\xi(2s)$，
其中 $\xi(s)=\pi^{-s/2}\Gamma(s/2)\zeta(s)$；于是 $E^{*}$ 的函数方程不过是 $\zeta$ 的函数方程的改写。
一般情形下 $A(s)$ 称为散射矩阵：它的极点给出 $L^{2}$ 中的剩余谱，即第 9 节的诸 $\eta_{j}$，
其行列式的相位描述散射的相移。
\par\vspace{0.15em}{\raggedleft 【来源：Iwaniec, 1995, Ch.~3 与 Ch.~13；Hejhal, Vol.~2, Ch.~6（A）】\par}
\end{annote}

$E(g,f_{s,k},\chi)$ 与 $c(s)$ 的解析延拓密切相关。Selberg 证明了：

\par\medskip\noindent\textbf{定理 41.}\quad\emph{$A_{k}(s)$ 与 $E(g,f_{s,k},\chi)$ 都有亚纯延拓，
且极点相同。它们在 $\operatorname{re}(s)>\tfrac12$ 内解析，只有实轴上 $(\tfrac12,1]$
内有限多个极点可能是例外。在直线 $\operatorname{re}(s)=\tfrac12$ 上它们全纯。
Eisenstein 级数满足函数方程}
\[
E(g,f_{s,k},\chi)=A_{k}(s)\,E(g,f_{1-s,k},\chi).
\]

\smallskip
我们只在 $\Gamma=\mathrm{SL}(2,\mathbb{Z})$ 的特殊情形证明它（定理 44）。
一般情形在文献中有各种叙述，见 Borel [5]、Cohen 与 Sarnak [9]、Colin de Verdi\`ere [8]、
Efrat [11]、Elstrodt [12]、Faddeev [13]、Harish-Chandra [19]、Hejhal [21, Chapter 6 与
Appendix F，文献讨论见 p.\ 225]、Kubota [29]、Jacquet [25]、Langlands [32]、
Lax 与 Phillips [34]、Moeglin 与 Waldspurger [36]、Osborne 与 Warner [37]、Venkov [48]、
Wong [53]。最一般的处理是 Langlands 那篇具有历史意义的工作 [32]，
以及 Moeglin 与 Waldspurger [36] 那部细致的现代专著。

这些文献都多少受惠于 Selberg，据 Hejhal [21]，他找到了三个证明。
一般而言，人们平行地证明 $A_{k}(s)$ 与 $E(g,f_{s,k},\chi)$ 的解析延拓。
基本原则是：算子的预解式解析延拓到其谱的补集。
这样用于 Kubota [29]、Venkov [48] 或 Langlands [32, Appendix IV] 等文献，
就得到 $A_{k}(s)$ 与 $E(g,f_{s,k},\chi)$ 到区域 $\operatorname{re}(s)>\tfrac12$、
$s\notin(\tfrac12,1)$ 的解析延拓。然而在此进路下，要得到整个平面上的亚纯延拓还有困难。

Selberg 与 Bernstein 各自独立地意识到，把这些困难绕开，可以把预解式原则与另一条原则结合起来。
后者的洞见是：若一个系数解析的线性方程组有\emph{唯一}解，
则其解延拓为亚纯函数，凡系数能延拓之处它就能延拓。
Selberg 关于此想法的已发表评论，见他的 G\"ottingen 讲义 [43] 1988 年撰写的导言。
除 Eisenstein 级数的解析延拓外，Bernstein 还给出了其他应用，
例如缠绕积分 $M(s)$ 及其 $p$ 进类似物 [2, 3] 的解析延拓，
他的方法已成为 $p$ 进群表示论的标准技巧。

在假定解析延拓成立的前提下，我们指出下面两个命题。

\par\medskip\noindent\textbf{命题 42.}\quad\emph{对一切 $k\equiv\epsilon$ 模 $2$，
函数 $E(g,s,f_{s,k})$ 构成一个同构于 $P_{s}^{\pm}$ 的 $(\mathfrak{g},K)$-模。}

\smallskip
\noindent\emph{证明.}\ 先设 $s>1$。此时 $E(g,s,f_{s,k})$ 由 $f_{s,k}$ 的左平移平均构成。
群及其李代数在函数上的表示通过右平移给出。两种运算彼此交换，故此情形的结论清楚。
由解析延拓，相应的恒等式保持。\hfill$\square$

\par\medskip\noindent\textbf{命题 43.}\quad\emph{若 $s$ 不是 $E(g,f_{s,k},\chi)$ 的极点，
则 $E-E_{0}$ 速降。}

\smallskip
\noindent\emph{证明梗概.}\ 设 $\phi\in\mathcal{H}$。函数
\[
\tilde E(g)=E(g,f_{s,k})-E_{0}(g,f_{s,k})
\]
关于 $\Gamma$ 不是自守的；但把 Bump [6] 中命题 3.2.3 的证明稍加修改（那里定义了核 $K_{0}$），
我们有
\[
T_{\phi}\tilde E(g)=\int_{\Gamma_{\infty}\backslash G}K_{0}(g,h)\,\tilde E(h,f_{s,k},\chi)\,dh .
\]
按定理 38 (ii) 的同一推理，可以选取 $\phi$，使 $T_{\phi}f_{s,k}$ 是 $f_{s,k}$ 的非零倍数。
现在 $\tilde E$ 落在 $V(k)$ 中，其中 $(\pi,V)$ 是同构于 $\pi_{s}^{\pm}$ 的函数空间，
故 $T_{\phi}\tilde E$ 是 $\tilde E$ 的同一个常数倍。速降性随即由核 $K_{0}(g,h)$
的相应性质推出。\hfill$\square$

由此可见，$E(g,s,f_{s})$ 在尖点附近的渐近性状由它的常数项决定。
常数项由两部分组成，一部分的量级是 $y^{s}$，另一部分是 $y^{1-s}$。
最小的增长显然出现在直线 $\operatorname{re}(s)=\tfrac12$ 上，那里量级是 $\sqrt{y}$。
在这条直线上，Eisenstein 级数\emph{几乎}平方可积，但并不平方可积。

算术情形（例如 $\Gamma=\mathrm{SL}(2,\mathbb{Z})$、$\chi=1$）与非算术情形有实质差别。
在算术情形可以归一化 Eisenstein 级数。这意味着存在 Eisenstein 级数的一个倍数，
它具有简单的函数方程，且只有有限多个极点。

例如设 $\Gamma=\mathrm{SL}(2,\mathbb{Z})$、$\chi=1$、$k=0$。则 Eisenstein 级数为
\[
E(z,s)=\sum_{\gamma\in\Gamma_{\infty}\backslash\mathrm{SL}(2,\mathbb{Z})}f_{s,0}(\gamma z),
\]
而\emph{归一化 Eisenstein 级数}是
\begin{equation}\tag{47}
E^{*}(z,s)=\pi^{-s}\Gamma(s)\zeta(2s)E(z,s).
\end{equation}

\par\medskip\noindent\textbf{定理 44.}\quad\emph{$E^{*}(z,s)$ 解析延拓到除 $s=0,1$ 以外的一切 $s$，
在这两点有单极点；$s=1$ 处的留数是常数函数 $\tfrac12$。
它满足函数方程 $E^{*}(z,s)=E^{*}(z,1-s)$。}

\smallskip
\noindent\emph{证明梗概.}\ 设 $z=x+iy\in\mathfrak{H}$、$t>0$，置
\[
\Theta(t)=\sum_{(m,n)\in\mathbb{Z}^{2}}e^{-\pi|mz+n|^{2}t/y}.
\]
由 $\Gamma$ 的 Euler 积分可得
\[
E^{*}(z,s)=\frac12\int_{0}^{\infty}(\Theta(t)-1)\,t^{s}\,\frac{dt}{t}.
\]
Poisson 求和公式蕴涵 $\Theta(t)=t^{-1}\Theta(t^{-1})$。由此得到
\[
E^{*}(z,s)=\frac12\int_{0}^{\infty}\Theta(t)\left(t^{s}+t^{1-s}\right)\frac{dt}{t}
-\frac{1}{2s}-\frac{1}{2-2s}.
\]
这个表达式给出解析延拓与函数方程。\hfill$\square$

定理 44 是定理 41 的特殊情形。它是算术情形的典型，却不是一般情形的典型。
\section{谱展开}

谱展开由 Roelcke [40] 得到，其前提是 Eisenstein 级数的解析延拓。文献中有许多叙述，
其中 Godement [17, 18] 的叙述很好且影响很大。本节讨论 $L^{2}(\Gamma\backslash\mathfrak{H})$
的谱展开，其中如第 6 节，群带一个尖点。我们假定 Eisenstein 级数的解析延拓成立（定理 41）。

设 $\phi$ 是 $C_{c}^{\infty}(G_{\infty}\backslash G)$ 中的 $K$-有限元素。
\emph{不完全 theta 级数}
\begin{equation}\tag{48}
\theta_{\phi}(g)=\sum_{\Gamma_{\infty}\backslash\Gamma}\phi(\gamma g)
\end{equation}
有点像 Eisenstein 级数，但它不是 $\Delta$ 的特征函数。
（Godement 的术语“不完全 theta 级数”似乎有点用词不当。）
若 $f\in C^{\infty}(\Gamma\backslash G)$，令
\[
f_{0}(g)=\int_{\Gamma_{\infty}\backslash G_{\infty}}f(ug)\,du
\]
为它的常数项，它属于 $C^{\infty}(G_{\infty}\backslash G)$。

\par\medskip\noindent\textbf{命题 45.}\quad\emph{不完全 theta 级数与常数项映射互为伴随，即}
\[
\int_{\Gamma\backslash G}\theta_{\phi}(g)\overline{f(g)}\,dg
=\int_{G_{\infty}\backslash G}\phi(g)\overline{f_{0}(g)}\,dg .
\]

\smallskip
\noindent\emph{证明.}\ 左端等于
\[
\begin{aligned}
\int_{\Gamma\backslash G}\sum_{\gamma\in\Gamma_{\infty}\backslash\Gamma}
\phi(\gamma g)\overline{f(g)}\,dg
&=\int_{\Gamma_{\infty}\backslash G}\phi(g)\overline{f(g)}\,dg\\
&\quad=\int_{G_{\infty}\backslash G}\int_{\Gamma_{\infty}\backslash G_{\infty}}
\phi(ug)\overline{f(ug)}\,du\,dg\\
&\quad=\int_{G_{\infty}\backslash G}\phi(g)\int_{\Gamma_{\infty}\backslash G_{\infty}}
\overline{f(ug)}\,du\,dg,
\end{aligned}
\]
这正是右端。\hfill$\square$

不完全 theta 级数在模 $\Gamma$ 意义下紧支，因而平方可积。
正因如此，它们比 Eisenstein 级数容易处理。

\par\medskip\noindent\textbf{命题 46.}\quad\emph{$L_{0}^{2}(\Gamma\backslash G)$ 是
$L^{2}(\Gamma\backslash G)$ 中由不完全 theta 级数张成的闭子空间的正交补。}

\smallskip
\noindent\emph{证明.}\ 由命题 45 立得，因为尖点谱的特征正是常数项为零。\hfill$\square$

本节余下部分作如下简化假设：特征 $\chi$ 平凡，且只考虑 $G$ 上右 $K$-不变的函数，
即可以看作 $\mathfrak{H}$ 上的函数。于是我们把 $E(g,f_{s,0},1)$ 简记为 $E(g,s)$，
或记作 $E(z,s)$，其中 $z=g(i)$。同时记 $f_{s}=f_{s,0}$。常数项写成
\begin{equation}\tag{49}
E_{0}(g,s)=f_{s}(g)+A(s)f_{1-s}(g).
\end{equation}
（见命题 40。）

利用定理 38，取 $\xi_{i}$ 为 $L_{0}^{2}(\Gamma\backslash G)$ 的一组基。又由定理 41，
Eisenstein 级数 $E(z,s)$ 可能在 $s_{j}\in(\tfrac12,1]$（$j=1,\dots,N$）处有有限多个极点，
这些也是 $A(s)$ 的极点。设 $\alpha_{j}$ 是 $A(s)$ 在 $s=s_{j}$ 处的留数，并令
\begin{equation}\tag{50}
\eta_{j}=\frac{1}{\sqrt{|\alpha_{j}|}}\operatorname{res}_{s=s_{j}}E(g,s).
\end{equation}

\par\medskip\noindent\textbf{命题 47.}\quad\emph{若 $\phi$ 是尖点形式，则}
\begin{equation}\tag{51}
\int_{\Gamma\backslash G}\phi(g)E(g,s)\,dg=0 .
\end{equation}

\smallskip
\noindent\emph{证明.}\ 先设 $\operatorname{re}(s)>1$。此时 (51) 为
\[
\begin{aligned}
\int_{\Gamma\backslash G}\sum_{\gamma\in\Gamma_{\infty}\backslash\Gamma}
\phi(g)f_{s}(\gamma g)\,dg
&=\int_{\Gamma_{\infty}\backslash G}\phi(g)f_{s}(g)\,dg\\
&=\int_{G_{\infty}\backslash G}\int_{\Gamma_{\infty}\backslash G_{\infty}}
\phi(ug)f_{s}(ug)\,du\,dg\\
&=\int_{G_{\infty}\backslash G}f_{s}(g)\int_{\Gamma_{\infty}\backslash G_{\infty}}
\phi(ug)\,du\,dg=0,
\end{aligned}
\]
因为 $\phi$ 是尖点形式。一般情形由解析延拓推出。\hfill$\square$

\par\medskip\noindent\textbf{命题 48.}\quad
(i) \emph{$\eta_{j}$ 的常数项是 $\sqrt{|\alpha_{j}|}\,f_{1-s_{j}}$。}

(ii) \emph{诸 $\eta_{j}$ 取实值、平方可积，且与尖点形式正交。}

\smallskip
\noindent\emph{证明.}\ (i) 由 (49) 立得，因为右端第一项没有极点而第二项有极点。
诸 $\eta_{j}$ 取实值，因为诸 $s_{j}$ 是实数。由命题 43，$\eta_{i}$ 在尖点附近与其常数项渐近，
故用 (i) 及 $s_{j}>\tfrac12$ 可知 $\eta_{i}$ 平方可积。
在命题 47 中取留数，即知诸 $\eta_{j}$ 与尖点形式正交。\hfill$\square$

\par\medskip\noindent\textbf{命题 49.}\quad\emph{设 $\phi\in C_{c}^{\infty}(G_{\infty}\backslash G/K)$，
定义}
\begin{equation}\tag{52}
\hat\phi(s)=\int_{G_{\infty}\backslash G}\phi(g)f_{1-s}(g)\,dg .
\end{equation}
\emph{它是 $s$ 的整函数，且对一切实 $\sigma$，$t\mapsto\hat\phi(\sigma+it)$ 属于 Schwartz 类。
对实 $\sigma$ 有}
\begin{equation}\tag{53}
\phi(g)=\frac{1}{2\pi i}\int_{\sigma-i\infty}^{\sigma+i\infty}\hat\phi(s)f_{s}(g)\,ds .
\end{equation}
\emph{（围道取实部为 $\sigma$ 的竖直线。）}

\smallskip
\noindent\emph{证明.}\ 我们有
\[
\hat\phi(s)=\int_{0}^{\infty}\phi\begin{pmatrix}y^{1/2}&\\&y^{-1/2}\end{pmatrix}
y^{-s}\,\frac{dy}{y}.
\]
$\phi$ 紧支且光滑，故 $t\mapsto\hat\phi(\sigma+it)$ 是 Schwartz 函数的 Fourier 变换，
因而本身属于 Schwartz 类。由 Mellin 反演公式
\begin{equation}\tag{54}
\phi\begin{pmatrix}y^{1/2}&\\&y^{-1/2}\end{pmatrix}
=\frac{1}{2\pi i}\int_{\sigma-i\infty}^{\sigma+i\infty}\hat\phi(s)\,y^{s}\,ds .
\end{equation}
这在角线上验证了 (53)；又两端都左 $G_{\infty}$-不变、右 $K$-不变，故一般情形随之成立。
\hfill$\square$

\par\medskip\noindent\textbf{命题 50.}\quad\emph{有}
\[
\langle\theta_{\phi},\eta_{j}\rangle=\sqrt{|\alpha_{j}|}\,\hat\phi(s_{j}).
\]

\smallskip
\noindent\emph{证明.}\ 由命题 45、命题 48 以及 $\hat\phi$ 的定义即得。\hfill$\square$

\par\medskip\noindent\textbf{定理 51.}\quad\emph{设 $\psi\in L^{2}(\Gamma\backslash\mathfrak{H})$。
再设 $\psi(g)f_{1/2}(g)$ 在 Siegel 集 $S$ 上可积。则}
\begin{equation}\tag{55}
\begin{aligned}
\psi(g)&=\sum_{j=0}^{N}\langle\psi,\eta_{j}\rangle\eta_{j}(g)
+\sum_{i=0}^{\infty}\langle\psi,\xi_{i}\rangle\xi_{i}(g)\\
&\quad+\frac{1}{4\pi}\int_{-\infty}^{\infty}
\left\langle\psi,E\left(-,\frac12+it\right)\right\rangle
E\left(g,\frac12+it\right)dt .
\end{aligned}
\end{equation}

\smallskip
关于 $\psi(g)f_{1/2}(g)$ 可积的假设，保证展开中出现的各内积收敛。
只要解释得当，此展开对一切 $f\in L^{2}(\Gamma\backslash\mathfrak{H})$ 都成立。
这类似于 $\mathbb{R}$ 的 Fourier 反演公式：它对一切 $L^{2}(\mathbb{R})$ 函数成立，
尽管作为积分来定义的 Fourier 变换只在 $L^{2}(\mathbb{R})\cap L^{1}(\mathbb{R})$ 上严格正确。

\smallskip
\noindent\emph{证明梗概.}\ 由命题 48 与命题 47，结论对尖点形式成立。
于是由命题 46，只需对不完全 theta 级数证明，即证明
\begin{equation}\tag{56}
\theta_{\phi}(g)=\sum_{j=1}^{N}\langle\theta_{\phi},\eta_{j}\rangle\eta_{j}(g)
+\frac{1}{4\pi}\int_{0}^{\infty}
\left\langle\theta_{\phi},E\left(-,\frac12+it\right)\right\rangle
E\left(g,\frac12+it\right)dt,
\end{equation}
其中 $\phi\in C_{c}^{\infty}(G_{\infty}\backslash\mathfrak{H})$。

事实上，两端都左 $G_{\infty}$-不变、右 $K$-不变，且由 (54) 在角线上相等，
故由 Iwasawa 分解即得。取 $\sigma>1$，对 $\gamma g$（$\gamma\in\Gamma_{\infty}\backslash\Gamma$）
求和，得到
\begin{equation}\tag{57}
\theta_{\phi}(g)=\frac{1}{2\pi i}\int_{\sigma-i\infty}^{\sigma+i\infty}
\hat\phi(s)\,E(g,s)\,ds .
\end{equation}
现在把积分路径向左移。$s=s_{j}$ 处的极点给出留数，由 (50) 与命题 50，
此留数等于 $\langle\theta_{\phi},\eta_{j}\rangle\eta_{j}$。于是
\begin{equation}\tag{58}
\theta_{\phi}(g)=\sum_{j=1}^{N}\langle\theta_{\phi},\eta_{j}\rangle\eta_{j}(g)
+\frac{1}{2\pi}\int_{-\infty}^{\infty}
\hat\phi\left(\frac12+it\right)E\left(g,\frac12+it\right)dt .
\end{equation}

现在来看 (56) 中的积分。由命题 45、(49) 与 (52)，被积函数是
\[
\begin{aligned}
&\frac{1}{4\pi}\int_{\Gamma\backslash G}
\theta_{\phi}(h)E\left(h,\frac12-it\right)dh\cdot E\left(g,\frac12+it\right)\\
&\quad=\frac{1}{4\pi}\left(\hat\phi\left(\frac12+it\right)E\left(g,\frac12+it\right)
+\hat\phi\left(\frac12-it\right)A\left(\frac12-it\right)E\left(g,\frac12+it\right)\right)\\
&\quad=\frac{1}{4\pi}\left(\hat\phi\left(\frac12+it\right)E\left(g,\frac12+it\right)
+\hat\phi\left(\frac12-it\right)E\left(g,\frac12-it\right)\right).
\end{aligned}
\]
对 $t$ 积分，合并重复的贡献，再用 (58)，即证得 (56)。\hfill$\square$

\begin{annote}
(55) 的三项分别对应剩余谱、离散谱与连续谱。$\eta_{j}$ 来自 $A(s)$ 在 $(\tfrac12,1]$ 上的极点，
只有有限多个；$\xi_{i}$ 是尖点形式；第三项是连续谱，被积函数 $E(g,\tfrac12+it)$ 本身不平方可积，
但按 Plancherel 测度 $\frac{1}{4\pi}dt$ 积分后仍给出 $L^{2}$ 中的展开。
这与第 8 节 $\mathbb{R}$ 上 Fourier 反演 $\phi(x)=\int\hat\phi(a)e^{2\pi iax}da$ 完全平行。
\par\vspace{0.15em}{\raggedleft 【来源：Selberg, J. Indian Math. Soc. 20 (1956)；Godement, Proceedings of Symposia in Pure Mathematics 9 (1966)；Iwaniec, 1995, Ch.~7（A）】\par}
\end{annote}

\section{提升与迹公式}

迹公式最有趣的应用之一，是自守形式的提升。方法是比较两个不同群上的两个迹公式，
由此得出一个群上的自守形式可以提升为另一个群上的自守形式。

Jacquet 与 Langlands [26] 给出了一个早期应用，即从除代数上的自守形式提升到 $\GL(2)$。
Gelbart 与 Jacquet [15] 注意到这一主题的一个变体，它大概是此类例子中最简单的一个，
因为此时两个迹公式都不涉及连续谱。改用阿代尔观点会比较方便，
但我们认为读过第 5 节的读者不难完成这一转换。

若 $D_{1}$ 与 $D_{2}$ 是域 $F$ 上的中心除代数，则 $D_{1}\otimes D_{2}\cong
\operatorname{Mat}_{k}(D_{3})$，其中 $D_{3}$ 与 $k$ 适当选取；
而 $D_{1},D_{2}\mapsto D_{3}$ 给出中心除代数的同构类之集 $B(F)$ 上的结合乘法。
于是 $B(F)$ 成为一个群，称为 \emph{Brauer 群}。

若 $D$ 是 $F$ 上的中心除代数，则 $D$ 的维数是平方数 $d^{2}$；又若 $E/F$ 是次数为 $d$、
可以嵌入 $D$ 的域扩张，则 $E\otimes D\cong\operatorname{Mat}_{d}(E)$。
于是除环是矩阵环的 Galois 扭形式。复合映射
\[
D\to E\otimes D\cong\operatorname{Mat}_{d}(E)\to E
\]
（最后一个映射取迹或行列式）取值于 $F$，给出\emph{约化迹}或\emph{约化范数}。

局部域或整体域 $F$ 的 Brauer 群有一个简单而漂亮的描述，与类域论的互反律相关。
见 Cassels 与 Fr\"ohlich [7] 中第 6 章（J.-P. Serre 的“局部类域论”）第 1 节，
以及第 7 章（J. Tate 的“整体类域论”）第 9 节。

设 $F$ 是整体域，$\A$ 是其阿代尔环，$D$ 是 $F$ 上次数为 $p^{2}$ 的中心除代数，
其中 $p$ 是素数。令 $Z\cong F^{\times}$ 为 $D^{\times}$ 的中心。
设 $S$ 是使 $D_{v}$ 为除环的有限位集。若 $v\notin S$，我们等同 $D_{v}=\operatorname{Mat}_{p}(F_{v})$。

以 $\mathcal{H}$ 记 $D_{\A}^{\times}=\prod_{v}D_{v}^{\times}$ 上这样一类函数之集：
它们是形如 $\prod_{v}\phi_{v}$ 的函数的有限线性组合，其中对每个 $v$，
$\phi_{v}:D_{v}^{\times}\to\mathbb{C}$ 光滑、模 $Z_{v}$ 紧支，当 $z_{v}\in Z_{v}$ 时满足
$\phi_{v}(z_{v}g_{v})=\phi_{v}(g_{v})$，且对 $F$ 的几乎所有位 $v$，
$\phi_{v}$ 与 $Z_{v}\operatorname{Mat}_{p}(\mathfrak{o}_{v})$ 的特征函数一致。
环 $\mathcal{H}$ 既含经典 Hecke 算子，也含上文第 7 节引入的积分算子。

$Z_{\A}D_{F}^{\times}\backslash D_{\A}^{\times}$ 紧。与 $\mathrm{SL}(2,\mathbb{R})$ 的情形一样，
$L^{2}(Z_{\A}D_{F}^{\times}\backslash D_{\A}^{\times})$ 对 $\phi\in\mathcal{H}$ 容许积分算子
$T_{\phi}$：
\[
(T_{\phi}f)(g)=\int_{Z_{\A}\backslash D_{\A}^{\times}}\phi(h)\,f(gh)\,dh .
\]
以 $\{\gamma\}$ 记 $D_{F}^{\times}$ 的共轭类的一组代表元。
用 $C_{\gamma}$ 记 $\gamma$ 在 $D_{F}^{\times}$ 中的中心化子。它是代数群，
故 $C_{\gamma}(\A)\subset D_{\A}^{\times}$ 表示它在 $\A$ 中的点。

\par\medskip\noindent\textbf{定理 52（Selberg 迹公式）.}\quad
\begin{equation}\tag{59}
\operatorname{tr}T_{\phi}=\sum_{\{\gamma\}}
\operatorname{vol}(C_{\gamma}\backslash C_{\gamma}(\A))
\int_{C_{\gamma}(\A)Z_{\A}\backslash D_{\A}^{\times}}\phi(g^{-1}\gamma g)\,dg .
\end{equation}

\smallskip
\noindent\emph{证明.}\ 命题 6 的证明原封不动地适用，故
\[
(T_{\phi}f)(g)=\int_{Z_{\A}D_{F}^{\times}\backslash D_{\A}^{\times}}K_{\phi}(g,h)f(h)\,dh,
\]
\[
K_{\phi}(g,h)=\sum_{\gamma\in D_{F}^{\times}/Z_{F}^{\times}}\phi(g^{-1}\gamma h).
\]
与 $\mathrm{SL}(2,\mathbb{R})$ 的情形一样，算子 $T_{\phi}$ 是 Hilbert--Schmidt 的；
再费些功夫可以证明它是迹类的。与定理 11 一样，
\[
\operatorname{tr}T_{\phi}=\int_{Z_{\A}D_{F}^{\times}\backslash D_{\A}^{\times}}
K_{\phi}(g,g)\,dg .
\]
于是 (59) 像定理 28 那样推出。\hfill$\square$

$D_{F}^{\times}$ 的共轭类容易描述。

\par\medskip\noindent\textbf{命题 53.}\quad\emph{若 $\alpha\in D_{F}^{\times}-Z_{F}$，
则 $F(\alpha)$ 是 $F$ 的 $p$ 次域扩张。$\alpha$ 与 $\beta$ 在 $D_{F}^{\times}$ 中共轭，
当且仅当存在域同构 $F(\alpha)\to F(\beta)$ 使 $\alpha\mapsto\beta$。
若 $F(\alpha)$ 是 $p$ 次域扩张，则 $F(\alpha)$ 可以嵌入 $D_{F}$，
当且仅当对一切 $v\in S$ 有 $[F_{v}(\alpha):F_{v}]=p$。}

\smallskip
\noindent\emph{证明.}\ $\alpha$ 与 $\beta$ 的共轭性由 Skolem--Noether 定理推出
（见 Herstein [22, p.\ 99]）。最后一个断言由 (i)$\Leftrightarrow$(ii) 推出
（见 Weil [51, Proposition VIII.5, p.\ 253]）。\hfill$\square$

迹公式可以用来证明许多情形下的函子性提升。

设 $E$ 是另一个次数为 $p^{2}$ 的除代数，并设 $E_{v}$ 为除环的位集与 $D_{v}$ 为除环的位集
$S$ 相同。若 $p=2$，这蕴涵 $D$ 与 $E$ 同构，但一般情形不然。
（这由整体类域论中 Brauer 群的计算推出，见 [7, Chapters 6 and 7]。）
因此我们要 $p>2$。

我们将证明 $D$ 上与 $E$ 上的两个自守形式空间同构。

设 $\pi=\otimes_{v}\pi_{v}$ 是 $L^{2}(Z_{\A}D_{F}^{\times}\backslash D_{\A}^{\times})$
的不可约成分。由于 $v\in S$ 时 $Z_{v}\backslash D_{v}^{\times}$ 紧，$\pi_{v}$ 在这些位有限维。
我们假定 $v\in S$ 时 $\pi_{v}$ 平凡。

若 $v\notin S$，则 $D_{v}\cong E_{v}\cong\operatorname{Mat}_{p}(F_{v})$。
于是 $v\notin S$ 时可以把 $\pi_{v}$ 与 $E_{v}$ 的一个不可约表示 $\pi_{v}'$ 等同，
而 $v\in S$ 时取 $\pi_{v}'=1$。令 $\pi'=\otimes\pi_{v}'$。
它是 $E_{\A}^{\times}$ 的不可约表示。

\par\medskip\noindent\textbf{定理 54.}\quad\emph{$\pi'$ 出现在
$L^{2}(Z_{\A}E_{F}^{\times}\backslash E_{\A}^{\times})$ 中。}

\smallskip
对应 $\pi\to\pi'$ 是 Langlands 意义上的自守形式函子性提升（见 Langlands [33]、Borel [4]）。

\smallskip
\noindent\emph{证明梗概.}\ 若 $v\in S$，则 $Z_{v}\backslash D_{v}^{\times}$ 紧，
故常函数 $\phi_{v}^{\circ}(g_{v})=1$ 属于 $C_{c}^{\infty}(F_{v})$。
以 $\mathcal{H}_{S}$ 记由诸函数 $\prod\phi_{v}$ 张成的 $\mathcal{H}$ 的子代数，
其中 $v\in S$ 时 $\phi_{v}=\phi_{v}^{\circ}$。它同构于 $E$ 上相应的 Hecke 环。
记此同构为 $\phi\to\phi'$。

由命题 53，$D_{F}^{\times}$ 与 $E_{F}^{\times}$ 中的非中心共轭类，
都与满足 $[F(\alpha):F]=p$ 且对一切 $v\in S$ 有 $[F_{v}(\alpha):F_{v}]=p$ 的域扩张元素
$\alpha$ 的 Galois 等价类之集一一对应。这一内蕴刻画表明，
我们可以把 $D_{F}$ 与 $E_{F}$ 的共轭类等同起来，再比较两个迹公式，得到
\begin{equation}\tag{60}
\operatorname{tr}T_{\phi}=\operatorname{tr}T_{\phi'}'.
\end{equation}
这并不像我们说的那么轻而易举，因为还须证明 (59) 右端的两个迹公式中的体积相同。

由 (60) 推出，$\mathcal{H}_{S}$ 在空间
\[
L^{2}\left(Z_{\A}D_{F}^{\times}\prod_{v\in S}D_{v}^{\times}\backslash D_{\A}^{\times}\right)
\quad\text{与}\quad
L^{2}\left(Z_{\A}E_{F}^{\times}\prod_{v\in S}E_{v}^{\times}\backslash E_{\A}^{\times}\right)
\]
上的两个表示同构，定理随即成立。\hfill$\square$

最后一步的底层事实是：环的两个表示由它们的迹刻画。例如，若 $R$ 是特征零的域上的代数，
$M_{1}$、$M_{2}$ 是有限维半单 $R$-模，且对每个 $\alpha\in R$，$\alpha$ 在 $M_{1}$ 与
$M_{2}$ 上诱导的自同态迹相同，则两模同构（见 Lang [31, Corollary 3.8, p.\ 650]）。
此陈述不能直接用于此处，但给出了思路。

余下部分取 $p=2$，并回顾 \emph{Jacquet--Langlands 对应}。
仍设 $D$ 如前。Jacquet--Langlands 对应是把自守表示从 $D^{\times}$ 提升到 $\GL(2,F)$。

对 $v\in S$ 有局部对应。$D_{v}^{\times}$ 模其中心紧，
故其不可约表示有限维。它们提升为 $\GL(2,F_{v})$ 的具有相同中心特征的不可约表示。
这一提升由 Jacquet 与 Langlands 用 theta 对应构造。事实上，$Z_{v}\backslash D_{v}^{\times}$
是正交群 $\mathrm{GO}(4)$ 的商，而 $\GL(2)$ 即 $\mathrm{GSp}(2)$，
theta 对应 $\mathrm{GO}(4)\leftrightarrow\mathrm{GSp}(2)$ 给出 Jacquet--Langlands 对应。
其像是平方可积表示（即超尖点表示与 Steinberg 表示）。

Jacquet 与 Langlands 先在其书第 14 节用逆定理，构造了从 $D^{\times}$ 上自守形式到
$\GL(2)$ 上自守形式的整体对应。为证明 $D^{\times}$ 上 $L$ 函数的函数方程，
他们用了 Godement--Jacquet 构造，因为此处没有 Hecke 积分可用。

最后，他们从迹公式的角度重新考虑了这一提升。这使他们能刻画提升的像。
他们写下了下述定理的证明梗概（后来由 Gelbart 与 Jacquet 完成）：

\par\medskip\noindent\textbf{定理 55.}\quad\emph{$\GL(2,\A)$ 的自守表示 $\pi$ 是
$D_{\A}^{\times}$ 的某个自守表示的提升，当且仅当对每个 $v\in S$，$\pi_{v}$ 平方可积。}

\smallskip
值得注意的是，他们对此事实的证明用到了此后三十年里变得重要的许多不同技巧：
$L$ 函数的 Hecke 与 Godement--Jacquet 积分构造、Weil 表示以及迹公式。

$\GL(2)$ 上的迹公式比除环上的难，因为有连续谱。我们通过证明定理 54 而不是定理 55
绕开了这一困难。

\begin{annote}
定理 55 就是 Jacquet--Langlands 对应：$D_{\A}^{\times}$ 的自守表示与 $\GL(2,\A)$ 中在每个
$v\in S$ 处为平方可积的自守表示一一对应，且对应保持 $L$ 函数与 $\epsilon$ 因子。
$p=2$ 时 $D$ 是四元数除代数，$S$ 是其分歧位集；局部对应由 theta 对应
$\mathrm{GO}(4)\leftrightarrow\mathrm{GSp}(2)$ 给出，其像正是超尖点表示与 Steinberg 表示。
整体对应亦可由比较两个迹公式得到，这正是 (60) 的用法。
\par\vspace{0.15em}{\raggedleft 【来源：Jacquet--Langlands, \emph{Automorphic Forms on GL(2)}, Springer LNM 114, 1970；Gelbart--Jacquet, Proc. Sympos. Pure Math. 33 (1979)（A）】\par}
\end{annote}

\section*{参考文献}
\addcontentsline{toc}{section}{参考文献}
\markright{参考文献}

\begin{enumerate}[label={[\arabic*]},leftmargin=2.8em,itemsep=0.35em,topsep=0.4em]
\item J. Arthur, The trace formula and Hecke operators, in K. E. Aubert, E. Bombieri
and D. Goldfeld, eds., \emph{Number Theory, Trace Formulas and Discrete Groups:
Symposium in Honor of Atle Selberg Oslo, Norway, July 14--21, 1987}, Academic Press,
New York, 1989, pp. 11--27.
\item I. Bernstein, \emph{Letter to Piatetski-Shapiro}, unpublished manuscript, 1986.
\item I. Bernstein, \emph{Meromorphic Continuation of Eisenstein Series},
unpublished manuscript, 1987.
\item A. Borel, Automorphic $L$-functions, in A. Borel and W. Casselman, eds.,
\emph{Automorphic $L$-functions: Automorphic Forms, Representations and
$L$-Functions, Part 2}, Proceedings of Symposia in Pure Mathematics 33, AMS,
Providence, RI, 1979, pp. 27--62.
\item A. Borel, \emph{Automorphic Forms on $\mathrm{SL}(2,\mathbb{R})$},
Cambridge University Press, Cambridge, UK, 1997.
\item D. Bump, \emph{Automorphic Forms and Representations}, Cambridge University
Press, Cambridge, UK, 1997.
\item J. Cassels and A. Fr\"ohlich, \emph{Algebraic Number Theory}, Academic Press,
New York, 1967.
\item Y. Colin de Verdi\`ere, Une nouvelle d\'emonstration du prolongement
m\'eromorphe des s\'eries d'Eisenstein, \emph{C. R. Acad. Sci. Paris S\'er. I Math.},
\textbf{293}(1981), 361--363.
\item P. Cohen and P. Sarnak, \emph{Discontinuous Groups and Harmonic Analysis},
unpublished lecture notes.
\item A. Connes, Noncommutative geometry and the Riemann zeta function, in
V. I. Arnold, ed., \emph{Mathematics: Frontiers and Perspectives}, AMS,
Providence, RI, 2002, 35--54.
\item I. Efrat, The Selberg trace formula for $PSL_{2}(\mathbb{R})^{n}$,
\emph{Mem. Amer. Math. Soc.}, \textbf{65}(1987), 359.
\item J. Elstrodt, Die Resolvente zum Eigenwertproblem der automorphen Formen in der
hyperbolischen Ebene I, II, III, \emph{Math. Ann.}, \textbf{203}(1973), 295--300;
\emph{Math. Z.}, \textbf{132}(1973), 99--134; \emph{Math. Ann.}, \textbf{208}(1974),
99--132.
\item L. Faddeev, The eigenfunction expansion of Laplace's operator on the fundamental
domain of a discrete group on the Loba\v{c}evski\u{\i} plane, Trudy Moskov. Mat.
Ob\v{s}\v{c}. \textbf{17}(1967) 323--350. \emph{Trans. Moscow Math. Soc.},
\textbf{17}(1967), 357--386.
\item S. Gelbart, \emph{Lectures on the Arthur--Selberg Trace Formula}, AMS,
Providence, RI, 1996.
\item S. Gelbart and H. Jacquet, Forms of $GL(2)$ from an analytic point of view,
in A. Borel and W. Casselman, eds., \emph{Automorphic Forms, Representations, and
$L$-functions, Part 1}, Proceedings of Symposia in Pure Mathematics 33, AMS,
Providence, RI, 1979, pp. 213--254.
\item I. Gelfand, M. Graev, and I. Piatetskii-Shapiro, \emph{Representation Theory
and Automorphic Functions}, Saunders, San Francisco, 1969; reprinted by Academic
Press, New York, 1990.
\item R. Godement, The decomposition of $L^{2}(G/\Gamma)$ for $\Gamma=SL(2,\mathbb{Z})$,
in \emph{Algebraic Groups and Discontinuous Subgroups}, in A. Borel and G. Mostow,
eds., \emph{Algebraic Groups and Discontinuous Subgroups}, Proceedings of Symposia in
Pure Mathematics 9, AMS, Providence, RI, 1966, 225--234.
\item R. Godement, The spectral decomposition of cusp forms, in A. Borel and
G. Mostow, eds., \emph{Algebraic Groups and Discontinuous Subgroups}, Proceedings of
Symposia in Pure Mathematics 9, AMS, Providence, RI, 1966, 225--234.
\item H. Harish-Chandra, \emph{Automorphic Forms on Semisimple Lie Groups}, Lecture
Notes in Mathematics 62, Springer-Verlag, New York, 1968.
\item D. Hejhal, \emph{The Selberg Trace Formula for $PSL(2,R)$}, Vol.\ 1, Lecture
Notes in Mathematics 548, Springer-Verlag, New York, 1976.
\item D. Hejhal, \emph{The Selberg Trace Formula for $PSL(2,R)$}, Vol.\ 2, Lecture
Notes in Mathematics 1001, Springer-Verlag, New York, 1983.
\item I. Herstein, \emph{Noncommutative Rings}, Carus Monograph 15, MAA, Washington,
DC, 1968.
\item H. Iwaniec, \emph{Introduction to the Spectral Theory of Automorphic Forms},
Biblioteca de la Revista Matem\'atica Iberoamericana, Madrid, 1995.
\item A. E. Ingham, \emph{Distribution of Prime Numbers}, Cambridge University Press,
Cambridge, UK, 1932.
\item H. Jacquet, Note on the analytic continuation of Eisenstein series: An appendix
to A. Knapp, ``Theoretical aspects of the trace formula for $GL(2)$,'' in
\emph{Representation Theory and Automorphic Forms (Edinburgh, 1996)}, Proceedings of
Symposia in Pure Mathematics 61, AMS, Providence, RI, 355--405.
\item H. Jacquet and R. Langlands, \emph{Automorphic Forms on $GL(2)$}, Lecture Notes
in Mathematics 114, Springer-Verlag, New York, 1970.
\item A. Knapp, \emph{Representation Theory of Semisimple Groups: An Overview Based
on Examples}, Princeton University Press, Princeton, NJ, 1986.
\item A. Knapp and E. Stein, Intertwining operators for semisimple groups,
\emph{Ann. Math.} (2), \textbf{93}(1971), 489--578.
\item T. Kubota, \emph{Elementary Theory of Eisenstein Series}, John Wiley, New York,
1973.
\item S. Lang, \emph{SL(2,$\mathbb{R}$)}, Addison--Wesley, Reading, MA, 1975.
\item S. Lang, \emph{Algebra}, 3rd ed., Addison--Wesley, Reading, MA, 1993.
\item R. Langlands, \emph{On the Functional Equations Satisfied by Eisenstein Series},
Lecture Notes in Mathematics 544, Springer-Verlag, New York, 1976.
\item R. Langlands, Problems in the theory of automorphic forms, in \emph{Lectures in
Modern Analysis and Applications III}, Lecture Notes in Mathematics 170,
Springer-Verlag, New York, 1970.
\item P. Lax and R. Phillips, \emph{Scattering Theory for Automorphic Functions},
Annals of Mathematics Studies 87, Princeton University Press, Princeton, NJ, 1976.
\item H. Maass, \"Uber eine neue Art von nichtanalytischen automorphen Funktionen und
die Bestimmung Dirichletscher Reihen durch Funktionalgleichungen, \emph{Math. Ann.},
\textbf{121}(1949), 141--183.
\item C. Moeglin and J.-L. Waldspurger, \emph{D\'ecomposition Spectrale et S\'eries
d'Eisenstein}, Birkh\"auser, Basel, 1993 (in French); \emph{Spectral Decomposition and
Eisenstein Series}, Cambridge University Press, Cambridge, UK, 1995 (in English).
\item M. S. Osborne and G. Warner, \emph{The Theory of Eisenstein Systems}, Academic
Press, New York, 1981.
\item S. Patterson, \emph{An Introduction to the Theory of the Riemann
Zeta-Function}, Cambridge University Press, Cambridge, UK, 1988.
\item B. Randol, Small eigenvalues of the Laplace operator on compact Riemann surfaces,
\emph{Bull. Amer. Math. Soc.}, \textbf{80}(1974), 996--1000.
\item W. Roelcke, \"Uber die Wellengleichung bei Grenzkreisgruppen erster Art,
\emph{S.-B. Heidelberger Akad. Wiss. Math.-Nat. Kl.} 1953/1955 (1956), 159--267.
\item A. Selberg, Harmonic analysis and discontinuous groups in weakly symmetric
Riemannian spaces with applications to Dirichlet series, \emph{J. Indian Math. Soc.
(N.S.)}, \textbf{20}(1956), 47--87.
\item A. Selberg, Discontinuous groups and harmonic analysis, in \emph{Proceedings of
the International Congress of Mathematicians (Stockholm, 1962)}, 1963, pp. 177--189.
\item A. Selberg, Harmonic analysis, in A. Selberg, \emph{Collected Works},
Springer-Verlag, New York, 1989, article 39; this is the annotated surviving portion
of the 1955 G\"ottingen lecture notes.
\item A. Selberg, On the estimation of Fourier coefficients of modular form, in
A. Whiteman, ed., \emph{Theory of Numbers}, Proceedings of Symposia in Pure
Mathematics 8, AMS, Providence, RI, 1965, 1--15.
\item T. Tamagawa, On Selberg's trace formula, \emph{J. Fac. Sci. Univ. Tokyo Sec. I},
\textbf{8}(1960), 363--386.
\item A. Terras, \emph{Harmonic Analysis on Symmetric Spaces and Applications I},
Springer-Verlag, New York, 1985.
\item V. Varadarajan, \emph{An Introduction to Harmonic Analysis on Semisimple Lie
Groups}, Cambridge University Press, Cambridge, UK, 1999.
\item A. Venkov, Spectral theory of automorphic functions, the Selberg zeta function
and some problems of analytic number theory and mathematical physics, \emph{Uspekhi
Mat. Nauk}, \textbf{34-3}(1979), 69--135; \emph{Russian Math. Surveys},
\textbf{34}(1979), 79--153.
\item A. Venkov, \emph{Spectral Theory of Automorphic Functions and Its Applications},
Kluwer, Dordrecht, The Netherlands, 1990.
\item A. Weil, Sur les ``formules explicites'' de la th\'eorie des nombres premiers,
\emph{Comm. S\'em. Math. Univ. Lund.}, volume dedicated to Marcel Riesz, 1952,
252--265.
\item A. Weil, \emph{Basic Number Theory}, Springer-Verlag, New York, 1967.
\item E. Whittaker and G. Watson, \emph{A Course of Modern Analysis}, 4th ed.,
Cambridge University Press, Cambridge, UK, 1927.
\item S.-T. Wong, The meromorphic continuation and functional equations of cuspidal
Eisenstein series for maximal cuspidal groups, \emph{Mem. Amer. Math. Soc.},
\textbf{83}(1990), 423.
\end{enumerate}
